% =====================================================================
%  R. Nielsen, OM BETINGELSERNE FOR EN KRAFTIG VILLIE.  Sex Forelæsninger.
%  Kjøbenhavn: Den Gyldendalske Boghandel (F. Hegel); trykt hos J. H. Schultz, 1874.  83 pp.  Antiqua.
%
%  TRANSCRIPTION -- GENERATED; DO NOT EDIT BY HAND.  Rebuild from the repo root:
%    python3 ebook-method/betingelserne-villie/draft.py --emit --force \
%      && python3 ebook-method/collate/check.py betingelserne-villie --italics --apply \
%      && python3 ebook-method/betingelserne-villie/draft.py --post
%
%  STATUS: COMPLETE (2026-10-12): title page and the whole text, pp. 1-83.  The book has no preface, contents or notes.
%
%  EDITION.  The text of record is the Google Books scan of the New York Public Library copy (barcode 33433081648804;
%  ~/bibliotek/Nielsen, Rasmus/betingelserne74.pdf, SCANS.tsv), 600 ppi bitonal, Google's OCR as text layer.
%  (betingelserne_for_en_kraftig_villie.pdf in the same folder is the same scan re-paged.)  The second printing,
%  »Andet Oplag« (Kjøbenhavn: J. H. Schubothe, 1884; KB scan betingelserne84.pdf, ABBYY layer), is a RESETTING with other
%  page breaks: it is not the text transcribed, but it served as a finding aid (below).  No e-book and no earlier
%  transcription exists (checked 2026-10-12: the repo, the web, KB).  An English translation, »The conditions of a
%  powerful will« (1882), exists and was not used.
%
%  ACCURACY (2026-10-12).  The words are those of Google's layer (draft.py), collated word by word with an independent
%  tesseract reading of the images (collate/check.py, witness 'ocr'): 623 disagreements.  78 were settled by check.py's
%  own rules, 213 by rules.py -- each rule grounded in this book's crops: the dashes the layer drops; o/ø and other
%  ø-faults of the layer where the ø-form is the corpus's word (the ø rule); the printer's space before ;:?! ; a layer
%  stop between a word and a lower-case word (a speck); tesseract's guillemets where the layer has none (the print sets
%  »...« and »...», both; the layer drops most of them and writes others as >> <<); the sheet signature.  The rest, 330,
%  were read at the image BLIND (sheets.py: crops with no readings shown, two planted controls a sheet, 58 of 58 controls
%  read right; six Sonnet readers, .parts/reads/r-b1-A.tsv, r-b2-B..F.tsv), and ~115 of those decided in the main session,
%  some at 300 dpi (zoom.py).  Then a FINDING AID: diff84.py compared the collated text with the 1884 resetting word by
%  word (236 differences) and every spot not explained by 1884's OCR or obvious revision -- 113 -- was read at the image
%  (zoom.py, z84-z90, zoom84-items.tsv).  It found what neither OCR could: ten spurious long dashes (tesseract reads a
%  wide space after a stop as „—“; the 'dash+' rule had taken them), two marks tesseract invented (pp. 20, 80), specks
%  read as stops by both OCRs, closing guillemets both missed, three words the layer split, a gutter loss (p. 40), and
%  six misprints.  Letterspacing (the print's only emphasis; no italics: slant scores 99% <= 1.02): candidates measured by
%  letter gaps (spacedsheets.py, threshold 0.095, 83 lines read blind; spaced.tsv, 56 runs); short spaced words off those
%  lines may be missed.  SAMPLED PAGE CHECK (pagecheck.py, seed 20261012): pp. 1, 27, 45 and the title page read whole
%  and blind, 929 words: 1 word error (the misprint „bebandle“ p. 27, which both OCRs had normalised), 0 punctuation,
%  0 markup -- 1.1 per 1000 words (95 % upper bound 5.1).  What this cannot see: a misprint both OCRs normalise where
%  the 1884 resetting has the correct word too -- the page check suggests about one in a thousand words.
%  PRINTED AS IS (each seen at the image; a % comment at the site): „netcp“ (p. 5), „Fortsaaelse“ (p. 11),
%  „Mennnesket“ (p. 12), „1“ for I (p. 17), „bebandle“ (p. 27), „Aabenbarlse“ (p. 54), „fortætte“ (p. 57),
%  „Følelsebedrag“ (p. 60); a missing stop after „Villie“ (p. 2), „Eenhed“ (p. 14), „Uleilighed“ (p. 82).
%  Supplied: p. 40 „tilvisse!“ (all but lost at the gutter in this copy) and „per[sonlige]“ (lost), read as the 1884
%  resetting has them.  The ø rule (Hans 2026-10-08) applied: „sløves“ (p. 82) and the layer's o-forms throughout.
%  Not transcribed: sheet signatures („R. Nielsen: Om Villien. N“, „N*“), library stamps and pencil, the copy's pen
%  strokes (one overdraws the dash on p. 36), the wavy rule under each lecture heading (\lecture draws a short rule).
%  Bold: one line, p. 49.  Quotation marks are as printed, in both directions.
%
%  PAGE MAP.  printed p = PDF page - 12 (p. 1 = PDF 13, p. 83 = PDF 95); title PDF 11, verso PDF 12 (stamps only).
%  Each lecture begins on a new page: I p. 1, II p. 12, III p. 23, IV p. 35, V p. 48, VI p. 65.
%
%  TOOLS (ebook-method/betingelserne-villie/, copied from om-aandsdannelse and udvalgte-stykker): draft.py, ocr.py,
%  rules.py, sheets.py + reads.py + READER-BRIEF.md, spacedsheets.py, diff84.py, zoom.py, pagecheck.py +
%  PAGECHECK-BRIEF.md, patches.py (PATCHES, POST), spaced.tsv, decisions.tsv.  Agent tokens: crop readers ~634k,
%  page reader ~88k.
%
% =====================================================================
\documentclass[12pt,a4paper,openany]{book}
\usepackage[utf8]{inputenc}
\usepackage[T1]{fontenc}
\usepackage{libertinus}
\usepackage{libertinust1math}
\usepackage[danish]{babel}
\usepackage[protrusion=true,expansion=true]{microtype}
\usepackage{setspace}
\usepackage[margin=1.3in]{geometry}
\usepackage{fancyhdr}
\usepackage{hyperref}
\hypersetup{pdftitle={Om Betingelserne for en kraftig Villie}, pdfauthor={Rasmus Nielsen}, hidelinks}
\pagestyle{fancy}
\fancyhf{}
\fancyhead[LE,RO]{\thepage}
\fancyhead[C]{\textit{R. Nielsen: Om Betingelserne for en kraftig Villie (1874)}}
\setstretch{1.3}
\usepackage{marginnote}
\renewcommand{\thefootnote}{\arabic{footnote})}
\newcommand{\tocline}[3]{\noindent\makebox[2.2em][r]{#1}\hspace{0.6em}#2\dotfill\ #3\par}
\newcommand{\opage}[1]{\marginnote{\footnotesize\textit{[#1]}}}
% a lecture opens a new page: heading, a short rule; the print's large initial is not reproduced
\newcommand{\lecture}[1]{\clearpage\begin{center}{\large #1}\\[0.4em]\rule{3em}{0.4pt}\end{center}\medskip\noindent}

\begin{document}

% --- front matter: PDF 11, title (to be read at the image).  The verso (PDF 12) carries only library stamps.
\begin{titlepage}
\begin{center}
\vspace*{3em}
{\Huge Om Betingelserne}\\[0.6em]
{\large for}\\[0.6em]
{\Huge en kraftig Villie.}\\[2.5em]
Sex Forelæsninger\\[1em]
af\\[1em]
{\large R. Nielsen.}\\[5em]
Kjøbenhavn.\\[0.5em]
Forlagt af den Gyldendalske Boghandel. (F. Hegel.)\\[0.5em]
{\small Trykt hos J. H. Schultz.}\\[0.5em]
{\small 1874.}
\end{center}
\end{titlepage}


% ===== TEXT (pp. 1--83) =====
% --- p. 1 ---
\setcounter{footnote}{0}%
\lecture{Første Forelæsning.}\opage{1}

Der er i Menneskenes Sind en ubetvingelig Uro, en
dunkel ustyrlig Drift, som jager efter Frihed. Hvad er da
Frihed? Et personligt Gode, altsaa mere end et udvortes
Vilkaar. Dersom Frihed blot var et udvortes Vilkaar, en
Tilstand, der kunde naaes ved at unddrage sig ydre Baand
og Tvang, da vilde Zigeuneren være friere end Statsborgeren,
Markens vilde Dyr og Fuglene under Himlen
friere end Mennesket. Men Frihed er væsenlig et indvortes
Gode, Frihed er vor Individualitets Selvstændighed, vor Personligheds
Grundpræg; Frihed er vort Adelsbrev, vort indre,
sande, oprindelige Selv.

Selvstændighed beroer paa Selvbestaaen, Selvbestaaen
paa Selvbestemmen, og Selvbestemmen er Frihed: jo mere
Selv der er i et Menneske, desto større er Friheden; jo
mindre og svagere Selvet er, desto mindre er Friheden.
Frihedens Grund er Villie, thi Villien er Selvets Kjerne.
Selvet er det Selvbestemmende, men Selvbestemmelsesevnen
er netop Villien; Villien er fri: jo kraftigere Villie, desto
større Frihed, jo svagere Villie, desto mindre Frihed.

Ere nu alle Mennesker oprindelig anlagte paa Frihed,
saa er jo hermed givet, at de alle uden Undtagelse maae
have en fri Villie. Men idet vi antage, at alle Mennesker
have fri Villie, saa maae vi dog tilstaae, at vi langtfra alle
besidde Villiens Frihed i lige Maal; Nogle have en stærk,
Andre en svag Villie: hvoraf kan det vel komme? Er
Villien noget Medfødt eller noget Selverhvervet eller maaskee
% --- p. 2 ---
\setcounter{footnote}{0}%
\opage{2}en Himlens Gave, en Gave af Aanden, den evigt frie? Vi
kunne nu gjerne sige, at den deels er noget Medfødt, deels
noget Selverhvervet, og dog i Et og Alt en Himlens Gave;
men saa kommer det da an paa, hvorledes dette nærmere
bliver at forstaae.

Vor første Forudsætning er, at Villien beroer paa
noget Medfødt. Til Villie hører naturligt Anlæg. Det
er en Kjendsgjerning, at der fødes Individualiteter med saa
fremtrædende Selvstændighed, at de fra Naturens Haand
udgaae som lutter Villie % PRINTED AS IS: no stop
Det er de fødte Herskere; Herskerpræget
med dets Overlegenhed er over dem, enten de
for Øvrigt befinde sig i høiere eller lavere Stilling. De,
der saae en Napoleon i hans første Ungdom, mødte det
samme Blik, de fine, læselige Skjæbnetræk, den samme Sluttethed
og jernfaste Villie hos den unge Elev fra Brienne som
siden hos den store Keiser. Andre derimod komme til Verden
med saa svage Selvstændighedsanlæg og forraade en saa
paafaldende Villieløshed, at de synes bestemte til at gaae
i Ledebaand hele Livet igjennem. Selv om de fødes i
Purpur og, saa at sige, med et Scepter i Haanden, er
Villiessvagheden dog saa stor, at de aldrig selv formaae
enten at fatte eller fastholde nogen afgjørende Beslutning.
De leve under et bestandigt Formynderskab: Yndlinge,
Tjenere, Raadgivere maae paatage sig at ville i deres Sted,
og de paatage sig det gjerne.

Istedenfor at gaae ud fra Grundsætningen om alle
Menneskers lige Adkomst til Frihedens Goder, skulde man
fra dette Synspunkt fristes til at troe, at Naturen havde
reist et uoverstigeligt Grændseskjel mellem Fribaarne og
Ufribaarne, havde forud bestemt Nogle til Selvstændighed,
Andre til Afhængighed, Nogle til Herrer, Andre til Slaver.
Vi tør imidlertid, hvad det Medfødte angaaer, ikke oversee,
at det ingenlunde forholder sig med Villiesevnen som med
andre Sjæleevner f. Ex. Hukommelsen. Naar En siger: jeg
har engang havt en meget tro Hukommelse, men desværre
næsten tabt den i en stor Sygdom, saa beklage vi ham,
og vor Beklagelse er uden Forbehold. Vil Nogen derimod
% --- p. 3 ---
\setcounter{footnote}{0}%
\opage{3}træde frem og sige: jeg har engang havt en meget kraftig
Villie, men ulykkeligviis mistet den i en Feber, saa beklage
vi ham vel ogsaa, men dog med et eget Forbehold. --- Man
kan miste Synet uden at miste Hørelsen, man kan til en
vis Grad tabe Hukommelsen uden at tabe Forstanden;
men vi kjende ingen Sygdom, hvori et Menneske kan tabe
Villien, og for Resten gjenvinde sin Helsen.

Den, der forsikkrer, at han har alle Sjæleevner ubeskaarne
paa Villien nær, som desværre gik tabt i en Koldfeber,
kan ikke være rigtig i Hovedet. Et besynderligt
Menneske! Han har Kundskaber, Dannelse, brillant Hukommelse
og som det synes ogsaa god Forstand, kan tale
med om mange Ting, Videnskab, Politik, Konst, Handelsaffairer,
uden at man mærker Noget; kun naar han kommer
paa det ene Punkt, det Punkt om Villien og Koldfeberen,
bliver han sær. Sligt kan hændes, naar vi tænke
paa Sindssygdomme. En ældre Kone, som i sin Familie
altid havde været anseet for sund og karsk, kom pludselig
til den Overbeviisning, at hun ikke kunde spise. Havde
hun da ingen Appetit? Jo. Var hun bange for at blive
forgivet? Nei. Hvad var der da i Veien? Ak, det var
en haard Skjæbne. Hun havde længe anet det; men nu
var det pludselig gaaet op for hende; hun var nemlig
ikke som andre Mennesker forsynet med indvendige Organer,
og hvad kunde det saa nytte, at hun spiste. Hun var,
forsikkrede hun, huul indeni, ganske huul, huul som en
Tromme, og saa græd hun. Paa lignende Maade kan et
Menneske falde paa den Fixidee, at han har tabt sin Villie,
eller slet ikke er blevet begavet med Villie, at der i hans
Indre er noget Huult, inderst inde et tomt Rum, et lille
mørkt Hul, just der, hvor Villien skulde sidde.

Det er en Fixidee; men idet vi føle os overbeviste
om, at det maa være en Fixidee, gjøre vi et Fremskridt i
Erkjendelse. Vi indsee nemlig, at Villien, trods det Medfødte
i Evnen, er anderledes afhængig af os selv end hvilkensomhelst
anden Evne. Jeg faaer ikke en stærk Hukommelse,
fordi jeg vil, eller en frugtbar Forstand, fordi jeg
% --- p. 4 ---
\setcounter{footnote}{0}%
\opage{4}vil, eller en dyb Følelse eller en rig Indbildningskraft,
fordi jeg vil; slige Fortrin beroe paa Naturbegavelse,
og herom gjælder det, at Ingen kan lægge en Alen til sin
Væxt. Med Villien er det en anden Sag; Villie faaer jeg
kun, idet jeg vil, ene og alene, fordi jeg vil. Et Menneske
kan have nok saa store Gaver og rige Evner: dersom
han ikke vil, saa faaer han ingen Villie. Den, der vil
halvt, faaer en halv Villie; men den, der kraftig vil, faaer
en kraftig Villie. Villien er fra dette Synspunkt sin egen
Aarsag: ved at ville frembringe vi Villie. Villien er altsaa
vor egen Frembringelse, en Frembringelse, hvorunder Selvet,
det Villende i os, frembringer sig selv. Det kan gives os at
ville, men Villien kan ikke gives os, vi maae erhverve den
selv: Villien er noget \emph{Selverhvervet}.

Hvorfor have saa ikke alle Mennesker en lige stærk
og kraftig Villie? Det maa, ifølge den sidst opstillede
Synsmaade, være deres egen Skyld. Og dog har ogsaa
denne Synsmaade sin Misviisning. Her melder sig En med
en Indvending, lad os høre ham. Han siger: Dersom Villie
alene beroede paa det at ville, saa vilde min Villie, det
være sagt uden al Beskedenhed, være en af de største og
stærkeste i Samtiden. Thi Villiens Selverhvervelse, Villie
frembragt ved Villie, en sig bestandig forøgende og forstærkende
Selvbestemmen, en mægtig, stor og kraftig Villie er
netop det Ideal, jeg har sværmet for fra min Ungdom af.
Men jeg er, just paa Grund af min Iver, kommet tilkort
med Idealet. Ikke har jeg hengivet mig til den sværmeriske
Indbildning, at et Menneske kan i det Udvortes Alt, hvad
han vil, at Villien kan flytte Bjerge. Jeg har holdt mig
til de mig forundte Evner; men netop i det Indvortes har
jeg seet den gabende Kløft mellem at \emph{ville} og at \emph{kunne}.
For at kunne bestemme Villien som min egen, maatte jeg
begynde med en omhyggelig og nøiagtig Udskillen af det,
som var mit, fra det, som ikke var mit. Men da jeg
havde udskilt Alt, som ikke var mit, hvormeget beholdt
jeg saa tilbage? Kun det at ville, den indre Bestemmen,
Villiestanken, den rene Villen, et Skyggevæsen uden Kraft.
% --- p. 5 ---
\setcounter{footnote}{0}%
\opage{5}Kun forsaavidt jeg er mig min Villie tydelig bevidst, tør
jeg kalde den min, og kun forsaavidt jeg tør kalde den
min, kan jeg selv bestemme den. Hvad der er mig ubevidst,
er ikke mit; Ingen kan bestemme, hvad han ikke
kjender. Det eneste i Villien, jeg er mig bevidst, er det,
at jeg vil. Med denne Bevidsthed kan jeg sidde hele Dage
og ville og ville uden at udrette det Allermindste. I
samme Øieblik jeg forsøger at foretage mig Noget, indløber
der strax en Mængde ueensartede Biomstændigheder, som
jeg ikke gjennemskuer, og hvorover jeg ikke er Herre. Jeg
kan sidde ved et rigtbedækket Bord, hungrig og tørstig:
dersom jeg ikke gjorde Andet end ville, saa sultede jeg
ihjel, og det var dog netcp % PRINTED AS IS (netop)
det, som jeg ikke vilde. Nu
sige Menneskene rigtignok, at Sligt er Spidsfindighed. »Ræk
Haanden ud og tag for dig; med andre Ord: brug dine
Kræfter, saa sulter du ikke ihjel «. Men jeg har en Mistanke
om, at Menneskene i slige Ting ikke altid vide,
hvad de selv sige. »Brug dine Kræfter!» det er jo netop
min Gaade. Er det Kræfternes Brug, der beroer paa Villien,
eller er det ikke snarere Villiens Brug, der beroer paa
Kræfterne? Naar vi vilkaarlig bevæge Arm og Haand, da
forekommer det os vel, at det er vor bevidste Villie, der
gjør det; men at det ikke er saaledes, kan den ringeste
Eftertanke lære os. Villien bevæger Nerverne, Nerverne
Musklerne, og Musklerne Armen, uden at vi paa
nogen Maade ere istand til at blive os bevidste, hvorledes
det egenlig gaaer til. Kun Saameget vide vi, at det ikke
er det i vor Villie, vi tør kalde vort Eget, som formaaer
at udrette Sligt. Bliver Armen lam, skjæres Bevægenerverne
over, saa kan den Paagjældende ville saa kraftig,
han vil: Armen bevæger sig ikke.

Dersom Kræfternes Brug ubetinget beroede paa vor
Villie, da maatte vor Villie jo selv have Kræfterne i sin
Magt. At erhverve sig Villie ved kraftig at ville, vilde da
bogstavelig være det Samme som at frembringe de legemlige
Kræfter, hvortil Villien trænger. Man maatte da ved
selverhvervet Villie kunne overvinde Sygdom og frembringe
% --- p. 6 ---
\setcounter{footnote}{0}%
\opage{6}Sundhed; men at dette er umuligt, tilstaaer Enhver. Det
er en Kjendsgjerning, at visse Muskler og Muskelbevægelser
ere Villien underlagte; men at det ingenlunde er Villien,
vor bevidste, selverhvervede Villie, som fra Begyndelsen
af har gjort sig til Herre over Musklerne, er indlysende.
At ville forklare Organernes Brug af Vanen er tankeløst.
Vane kommer ved Gjentagelse; Villien, hedder det, gjentager
Brugen af Organerne saalænge, indtil den tilsidst
bliver Herre over Bevægelserne. Saa forklarer man Kræfternes
Brug ved at sige: Brugen gjentager sig; en maadelig
Forklaring, naar der netop spørges om Brugens Mulighed.
Det er jo, som vilde man forklare et Mirakel ved
at sige: Miraklet gjentager sig.

Den Sygdom, man kalder St. Veits Dands, giver
iøinefaldende Exempler paa Misforhold mellem Villien og
Bevægelserne. Lemmerne bevæge sig uden at adspørge
Villien. --- I Talens Strøm synes det, som om Ordene i
deres Forbindelse ligefrem bleve styrede af vor bevidste
Villie, og dog er det ikke altid let at finde just det Ord,
man søger. I Tilfælde af Hjernesvækkelse kan Skilsmissen
mellem det, den Syge egenlig vilde sige, og det
han virkelig siger, endog være saa stor, at han f. Ex.
siger Kjøkken, naar han mener Brød, og Skuffe, naar han
mener Penge. Villien sigter paa det Rette, men de bevægende
Kræfter unddrage sig den rette Brug.

Det er med Villien som med en Herre, der udretter
Alt ved Tjenere, og desaarsag maa bære Ansvaret
for Alt, hvad Tjenerne gjøre. Er det noget Godt, saa
roser man, ikke Tjenerne, men Herren, thi det var jo ham,
der vilde det; er det noget Ondt, saa dadler man, ikke
Tjenerne, men Herren, thi det var jo ham, der vilde det.
Dette er altsammen i sin Orden, saalænge Tjenerne ere,
som villieløse Slaver, punktlig lydige. Men begynde de
at gaae paa deres egen Haand og blive ulydige eller dog
kun halvt lydige, saa kan Herren ikke længer tage Skylden
paa sig. Og dog har en Herre i Reglen større Hævd
over sine Tjenere, end Villien over sine Kræfter. Det
% --- p. 7 ---
\setcounter{footnote}{0}%
\opage{7}staaer i Herrens Magt at give en ulydig Tjener Afsked
og antage en ny Tjener; men det staaer ikke i Villiens
Magt at give de gamle ulydige Nerver Afsked og antage
nye i deres Sted.

Brug dine Kræfter! det er som vilde en Aladdin
sige: gnid paa Lampen! Villiens Brug af Kræfterne er
vel ikke eventyrlig, thi den er bunden til Lovene; men i
sin sidste Grund er den desuagtet lige saa uforklarlig som
Aandens Komme, naar Lampen gnides. Villien vil, og
Kræfterne røre sig; der gnides paa Lampen, og Aanden
kommer: det første Phænomen er gaadefuldt som det
sidste.

Vi have her betragtet to indbyrdes modsatte Forestillingsmaader,
af hvilke den første udelukkende gjør
Villien til noget \emph{Medfødt}, medens den sidste ligesaa
udelukkende gjør den til noget \emph{Selverhvervet}. At de
begge ende i det Urimelige, kommer af, at de begge, hver
paa sin Maade, forvexle Villien selv med dens Betingelser.
Det Medfødte er endnu ikke Villie; det er kun Villiesanlæg,
og forsaavidt en \emph{Betingelse}. Det Selverhvervede er ikke
Villiens Grund; Selverhvervelsen forudsætter netop det, som
skal erhverves. For at frembringe Villie, maa man ville, men
for at \emph{kunne} ville, maa man jo allerede have Villie, for
kraftig at ville, maa man have en kraftig Villie. Den,
som slet ingen Villie har at begynde med, kan ikke ville,
og den, som ikke forud har en kraftig Villie, kan ikke
ville kraftig. Saa sees da heraf, at den Selverhvervelse,
hvorunder Villien bliver til og voxer, ikke kan være Villiens
inderste Grund og første Aarsag, men kun en Betingelse.


Nu kunde man fristes til at slaae de to Betingelser
sammen og sige: Villien er deels noget Medfødt, deels
noget Erhvervet; ved at tage det Medfødte og det Erhvervede
under Eet, faae vi baade Villiens Selvbestemmen
og dens Grund, altsaa: den hele Villie. Det er en Umulighed;
de to Betingelser ere altfor ueensartede til, at de
saadan uden videre skulde lade sig sammensmelte. Der
% --- p. 8 ---
\setcounter{footnote}{0}%
\opage{8}er i Villiens Væsen Noget,  der beroer paa Naturen, og
Noget, der beroer paa os. Men det, der beroer paa Naturen,
kan umulig beroe paa os; thi, hvis saa var, maatte
vi jo sige, at det Hele beroer paa os. Omvendt kan det,
som beroer paa os, da heller ikke siges at beroe paa Naturen;
thi, hvis saa var, maatte jo det Hele komme til
at beroe paa Naturen. Nei, det, som skal kunne forene
de to ueensartede Betingelser, maa nødvendig være et
tredie Høiere, en Villiesmagt, som paa en Gang er i Naturen
og dog over Naturen, i os og dog over os: denne
Villiesmagt er \emph{Aanden}.

Idet vor frie Villie opfattes som en Aandens Gave, bliver
den i allerhøieste Forstand en Opgave, og Opgavens Løsning
religiøs. Men netop ved det Religiøse begynde Vanskelighederne
forfra, thi her begynde Mistydningerne for Alvor.

Det Religiøse har mange Fjender. Fromme Mænd
klage bitterlig over Nutidens Ugudelighed og tiltagende
Vantro; de have vel Grund dertil, naar de vende Blikket udad
og see det stadig voxende Tal af udvortes Fjender. Men
Fjenden er i det Indvortes! »Ja«, svare de, »vor Arvefjende
er indvortes, det er Djævelen, Antichristen, Dyret
med de mange Horn, hvis Tegn vi snart faae at see --- i
det Udvortes«. Men Fjenden er i det Indvortes, i det Religiøse
selv, en Hær af snigende Misforstaaelser. De første
farlige Fjender, et troende Sind har at værge sig imod,
ere Misfostre af dets egen opulmende Indbildning. Slige
Misfostre vise sig under mange Forklædninger; de melde
sig ikke altid med Sjæleangst, lyssky Forsagthed og Tvivl, de
melde sig ligesaa ofte med svullen Tro, falsk Sikkerhed
og en vis tykblodig Selvtillid, som hindrer Aandens
Udvikling. De, som ville bygge paa Religionens hellige
Grund, maae bære sig ad som Israeliterne, da de byggede
paa det andet Tempel, den Helligdom, som i Tidens Fylde
skulde modtage Verdens Frelser; de byggede med Værget
i Haanden. Kun er Forskjellen den, at Fjenderne af hiint
udvortes Tempel vare udvortes; men Fjenderne af Aandens
Helligdom ere væsenlig indvortes.

% --- p. 9 ---
\setcounter{footnote}{0}%
\opage{9}Ogsaa Reflexioner ere farlige, naar Forstanden, medens
den endnu er usikker over Villielivets Forudsætninger,
begynder at raisonnere. Her er En, som klager, at han
trods al sin Eftertanke ikke er istand til at komme paa
det Rene med, hvorledes det i Troen egenlig forholder sig
med Aanden og Villien. Jeg hører jo nok, siger han, en
heel Deel salvelsesfulde Forsikkringer om, at det er Gud
den Helligaand, der gjenføder Mennesket, skaber i os et
nyt Menneske, et nyt Sind, en ny Villie o. s. v. Naar
der tales saadan i al Almindelighed derom, især naar der
tales i Pathos, faaer jeg en Fornemmelse, som om jeg ogsaa
forstod det, at sige: i al Almindelighed. Men naar jeg
saa gaaer i mig selv og ret forsøger at gjøre mig Rede for,
hvorledes det da egenlig forholder sig med Aand og Villie,
saa løber det rundt for mig. De Opvakte tilraabe mig:
»Du skal ikke gruble, men troe; naar Du bare vil troe, saa
gaaer det som af sig selv!« Jeg vil gjerne troe, men hvormed
skal jeg praktisk begynde, fra hvem skal Begyndelsen udgaae,
fra Aanden eller mig selv? Er det mig, der skal
gjøre det første Skridt, saa nødes jeg da til at begynde
med den Villie, jeg har; men den Villie, jeg har, er jo
netop den gamle Villie, en Villie, som ikke duer, en Villie
udenfor Aanden, altsaa en aandløs Villie. Hvorledes kan
et Menneske komme til Aanden ved Hjælp af en aandløs
Villie? Saa hedder det: »Du skal bede om Aanden!« Men
saalænge jeg endnu ikke har faaet Aanden, er jo min Bøn
aandløs! Det kan altsaa ikke være mig, der skal begynde;
det maa være Aanden. Staaer det nu fast, at Begyndelsen
udgaaer, ikke fra mig, men fra Aanden: saa har jeg intet
Andet at gjøre end at sidde stille hen og vente paa Aanden.
Skal jeg da blive Kvæker? Man svarer mig: »Du
har jo Aanden, saa sandt Du er døbt. Det kommer alene
an paa Dig selv, om Du vil ihukomme din Daabspagt og
lade Aanden virke.« --- Lade Aanden virke! Men saa
er her jo dog Noget, der udelukkende beroer paa mig,
ikke paa Aanden. Hvad vil det sige, at ihukomme sin
Daabspagt? Det kan dog umulig være en blot og bar
% --- p. 10 ---
\setcounter{footnote}{0}%
\opage{10}Hukommelsessag; her maa være Noget, som udtrykkelig
beroer paa Sind og Villie, mit eget Sind, min egen
Villie. For at Aanden skal kunne virke i mig, maa
jeg altsaa, som ovenfor sagt, ved min, endnu aandløse,
Villie hjælpe den til at virke. Saa griber man
da til den sidste Forklaring og siger: »Du skal ikke
belure Aanden, ikke formaste Dig til at ville forstaae, hvorledes
den virker. Den virker i det Skjulte som en ny Natur;
Du skal erfare dens Virksomhed som en Kjendsgjerning.
Du kan erfare, at Planten voxer, men ikke gjennemskue
Naturens Hemmelighed og begribe, hvorledes. Det nye Liv
er en ny Plante, det voxer en stille Væxt: vi voxe, vi
groe, det er en Kjendsgjerning.»

Saa føres der da megen opbyggelig Tale om Livets aandelige
Væxt, hvorledes det groer, selv naar vi sove. Jeg kan
desværre ikke faae Øie paa Kjendsgjerningen. Bygget groer
og Havren groer; det groer i den Frugtsommelige, og
Barnet fødes og voxer til Mand, og der vaagner i Manden
en Drift, som Barnet ikke kjendte; det er Kjendsgjerninger.
Men med Livets Væxt i de Troende er det en egen Sag.
Jeg har i mange Aar hørt dem sige, at det groer i dem,
at de bestandig voxe; jeg har taget dem paa Ordet og
saa nøiagtig som muligt taget Maal af dem, for dog tilsidst
at opbygges ved Kjendsgjerningen; men hvad har jeg
erfaret? I det Udvortes har jeg vel kunnet see, at ogsaa
de Troende forandre sig med Aarene, at Nogle af dem
blive, om ikke altid høiere, saa dog ikke sjældent endeel
førere; men hvad Charakteerudvikling angaaer, har jeg ikke
været istand til at opdage nogen Fremgang. Det forekommer
mig, at Gjennemsnittet staaer akkurat ved det
samme Maal iaar som for ti, ja for tyve Aar siden: skikkelige
Mennesker, som --- paa Talemaaderne nær --- have en
paafaldende Lighed med andre skikkelige Mennesker. Det med
Væxten maa da være en stor Hemmelighed. Min Hemmelighed
er, at det ikke vil groe i mig. Jeg fik ogsaa i
sin Tid en Opvækkelse og haabede, at det skulde groe;
% --- p. 11 ---
\setcounter{footnote}{0}%
\opage{11}men desværre, det groer ikke: jeg er som Grantræet, der
engang havde været Juletræ, --- sat i Væxten.

Denne Prøve paa tilraisonnerede Misforstaaelser indeholder
en kraftig Paamindelse om, at Fjenden er i det
Indvortes.

Hvad sigter saa alt dette til? En Fortsaaelse % PRINTED AS IS (Forstaaelse)
af,
hvad Villie er. Vi begyndte med at antage Villien for
noget Medfødt; men Antagelsen førte til det Urimelige.
Vi fortsatte med at opfatte Villien som noget Selverhvervet,
men Opfattelsen førte til det Urimelige. Vi endte med at
henføre Villien til Aanden; men idet Opgaven blev religiøs,
kom Misforstaaelserne. Saa skulde dette da indlede
et Forsøg paa at eftervise Villiens Udspring og \emph{klare Betingelserne
for en kraftig Villie}.

% --- p. 12 ---
\setcounter{footnote}{0}%
\lecture{Anden Forelæsning.}\opage{12}

Efter en gammel Inddeling siges Mennesket at bestaae
af tre Hoveddele: \emph{Legeme}, Sjæl og \emph{Aand}. Denne Tredeling
maa dog ikke forstaaes, som om Mennnesket % PRINTED AS IS
var
sammensat af tre Stykker, og det Hele fremkommet ved,
at Sjælen var sat ind i Legemet, og Aanden ind i Sjælen,
at det Ydre skulde omgive det Indre som den ene Æske
kan omgive den anden. Dersom det Indre kun var indæsket
i det Ydre, vilde jo Eenheden mangle, den organiske
Eenhed, der gjør os til levende Væsener.

Et af de stærkeste Vidnesbyrd for, at det menneskelige
Legeme, trods Mangfoldigheden og Forskjelligheden af dets
Dele, ikke er, hvad man kalder en Maskine, et sammenstykket
Hele, men besjæles af Liv og Selveenhed, er Følelsen
af Smerte. Der kan være Smerte i de forskjelligste Dele
af Legemet, men Smerten som Smerte vedkommer egenlig
ikke Delene for sig, ikke engang Nerverne; det er Sjælen,
der føler. Den, der, for at holde os til Hverdagserfaringer,
har f. Ex. Tandpine, siger ikke: det maa Tænderne om, det
er en Uro i Nerverne, som ikke vedkommer mig; nei, han
maa bekjende, at det vedkommer ham, at det er ham selv,
der maa føle Smerten, den Smerte, som er i Tænderne.

Var Mennesket et saadant sammenstykket Hele, at
hver Deel havde sit eget Liv og sin egen Følelse for sig,
da maatte det være det fuldt bevidste Selv let at underkaste
sit Legeme Operationer; men for at underkaste sig
% --- p. 13 ---
\setcounter{footnote}{0}%
\opage{13}smertelige Operationer og dog, saavidt muligt, unddrage sig
Smerten, bliver man jo nuomstunder chloroformeret. Ved
denne Selvbedøvelse skeer det, at Legemets enkelte Dele,
de enkelte Organer, blive som Stykker, der lide hvert
for sig.

Følelsen, Selvfornemmelsen især, voxer, udvikles, forfines
med Civilisationen; men det seer næsten ud, som om
Villieskraften tog af. Da Mucius Scævola stak sin Haand
i Ilden og lod den brænde, medens han talede til Porsenna,
da var han ganske vist ikke chloroformeret. Vel seer det
ud, som var hans Haand kun et Stykke Legeme for sig,
som kom det ikke ham ved, men var alene Haandens Sag,
hvorvidt den kunde taale Ilden eller ikke; og dog var det
ingen Udstykning. Han følte, hvad der foregik, han følte
Smerten; men der var i hans Selv et Noget, som var
stærkere end Følelsen, hævet over Smerten, det var hans
besluttede Villie. Mucius Scævola var en Mand med kraftig
Villie. Er det saa maaskee kun Manden, der i Forhold til
Smerten kan have en saadan kraftig, heroisk Villie? Nei!
ogsaa Kvinden. Historien veed at tale derom i Nordboernes
Liv og i Romernes. Da Mrc. Brutus, Cassius og Casca
havde sammensvoret sig for at myrde Cæsar og frelse Republiken,
bar Brutus tungt paa Hemmeligheden og turde ikke
betroe den til sin Hustru: hun var jo dog kun en Kvinde.
Porcia anede, hvad der foregik, men for at gjøre sig værdig
til en Indvielse og vise, at hun, en Kvinde, kunde staae ved
Mandens Side selv i et Rædslernes Øieblik, gav hun sig
et dybt og smerteligt Stik i Foden. Saa traadte hun frem,
let og med sikkre Fjed, skjøndt Smerten brændte, og nu
turde Brutus betroe hende Alt. Porcia var en Kvinde med
Villie, en stor og fast og kraftig Villie.

Det er da klart af Villien saavel som af Følelsen, at
vort menneskelige Væsen ikke er et sammenstykket Hele,
men en Livseenhed, en sig bestemmende Selveenhed, en
Følelses- og Villieseenhed, et Jeg.

Den sunde Sands har paa sin Viis hævet sig over
Sammenstykningen ved at løfte Mennesket op over Trede%
% --- p. 14 ---
\setcounter{footnote}{0}%
\opage{14}lingens tre Dele; hvert Menneske kalder sig et Jeg og siger:
Jeg har et Legeme, jeg har en Sjæl, jeg har Aand.

Her indtræder nu et besynderligt Forhold mellem:
at have og at være. Hvad vor Besiddelse af udvortes
Ting angaaer, falder det unaturligt at gjøre det, vi have,
til Eet med det, vi ere. Fordi En har en Herregaard,
siger han dog ikke: jeg er, hvad jeg har, en Herregaard,
saalidt som den, der har en Stok, vil sige: jeg er en Stok.
Rigtignok skal det i Amerika i nogle Kredse være Skik, at
Værten præsenterer to Herrer i Selskabet for hinanden kort
kategorisk som: to Millioner Dollars! --- fire Millioner! men,
ihvorvel en saadan Talbestemmelse har den Fordeel, at den
exact angiver Inclinationsvinklen for det forskjellige Nik, hvormed
de To bør komme hinanden imøde: saa kan det dog ikke
negtes, at en saa ligefrem Sammensmeltning af det at have
med det at være maa forekomme os noget vel materialistisk.

Anderledes, naar Talen er om vort Væsens Bestanddele
og vore væsenlige Egenskaber. I denne Sammenhæng kan
det at have paa forskjellig Maade falde sammen med det
at være. Vi have et Legeme, vi ere Legemer, vort Væsen
er legemligt; vi have en Sjæl, vi ere Sjæle, vort Væsen er
sjæleligt; vi have Aand, vi ere Aander, vort Væsen er
aandeligt. I denne Sammensmeltning af, hvad vi have og
hvad vi ere, fastholde vi vort Væsens Eenhed % PRINTED AS IS: no stop
Vort Selv,
vort Jeg, det Jeg, vi have, og det Jeg, vi ere, er fra første
Spire af legemligt-sjæleligt-aandeligt. Der er i Mennesket
kun een Natur; men denne Natur har tre Hovedsider, og
den hele Natur er igjen i enhver af sine Hovedsider. Der
er i vor Grundnatur tre Naturer: en \emph{sandselig}, svarende
til Legemligheden og den Side af Sjælen, som gaaer ind
under Sandsning og sandselige Tilskyndelser; en \emph{fornuftig},
svarende til det Ulegemlige, altsaa til den Side af Sjælen,
som gaaer op i Ideeanskuelser og Tænkning af det Almene;
en \emph{aandelig}, svarende til det oprindelige Selv i evig Frihed,
til Aanden.

Denne Tredeling er afgjørende for Villien og dens
Betingelser, thi heraf følger, at ogsaa Villien er tredeelt.
% --- p. 15 ---
\setcounter{footnote}{0}%
\opage{15}Vore mange Bryderier med Friheden, alle vore Anstrengelser
med at faae os selv i vor Magt og arbeide os op til Selvstændighed
hidrøre fra, at vi oprindelig have tre Villier:
en \emph{Driftvillie}, svarende til vor sandselige Natur og det
i Villien Medfødte, en \emph{Fornuftvillie}, svarende til vor
Erkjendelse af de almeengyldige Verdenslove og det i Villien
Selverhvervede, og endelig en \emph{Aandsvillie}, en Villie, som,
naar den forstaaer sig selv og erkjender sit Udspring,
har sin Selvstændighed i ubetinget Afhængighed af Aanden
selv, og veed sig bevæget af Guds Villie.

Der er altsaa tre Trin, hvorpaa Mennesket kan staae
med sin tresidige Villie; men det er klart, at vi maae staae
enten paa det første eller paa det andet eller paa det
tredie af disse Trin, naar vi skulle have en udeelt, fast og
kraftig Villie.

\emph{Naturmennesket}, hvorved vi ikke blot forstaae
Mennesket som det er i raa, naturlig Tilstand, men det
til sin medfødte Naturgrund umiddelbart bundne Menneske,
staaer paa det første Trin. Det er med sin individuelle
Begavelse, sin Fornuft og Genialitet som af een Støbning.
Naturmennesket har, uden Hensyn til, hvorvidt dets Civilisering
er naaet, kun een herskende Villie, een Grundvillie,
og det er \emph{Driftvillien}. Heltelivet i det gamle Norden
og det hos Grækerne blomstrende Kulturliv lærer os, hvad
dette vil sige.

\emph{Fornuftmennesket}, hvorved vi udelukkende tænke
paa Saadanne, hvis Villie bevæges af det for Tanken Nødvendige
og i sin væsenlige Nødvendighed Almeengyldige,
staaer paa det andet Trin. Tænkeren, den Vise, der kun
vil, hvad han maa ville, har i sin Overeensstemmelse med
sig selv een Villie, een Grundvillie og det er: \emph{Fornuftvillien}.
Oldtidens store Charakterer, en Sokrates, Stoikerne
o. Fl., afgive Exempler.

\emph{Aandsmennesket}, i Forbilledet selv det eneste
sande Menneske, Frihedens Aabenbarer, der vil frigjøre Alle,
staaer udelukkende paa det tredie Trin. Her er i Villieseenhed
to Villier, en guddommelig og en fuldstændig
% --- p. 16 ---
\setcounter{footnote}{0}%
\opage{16}menneskelig; men de to Villier ere ikke udstykkede, ikke
splidagtige; den guddommelige er menneskelig og den
menneskelige guddommelig; begges Eenhed er en absolut
personlig Villieseenhed, det er \emph{Aandsvillien}.

Men vi, det nittende Aarhundredes Børn, Natur-, Kultur- og
Aandsmennesker af mange Kuld, hvor staae nu
egenlig vi? Ikke paa noget enkelt af de tre Trin; vi staae
paa alle Trin, alle mulige Mellemtrin. Det kommer af, at
vi ere saa alsidige. At staae paa eet Trin, at blive staaende
og staae fast, er jo eensidigt, og Eensidighed er den Udyd,
man af alle mest maa frygte. Saa ere vi da alsidige, det
vil sige: vi have Noget af en Driftvillie, Noget af en Fornuftvillie
og Lidt af en Aandsvillie i forskjellig Blanding.
Denne vor Alsidighed bevirker, at Villielivet trods al dets
Rørighed gaaer noget trangt. Vi maae tjene ikke blot to,
men tre Herrer; thi enhver af de tre Villier i os er en
Stormagt, og Stormagter ere altid herskesyge. Kan Ingen
af dem blive Overherre, saa vække de Splid; de ville svække
os, udstykke vort Væsen og dele os, som Stormagterne i
sin Tid deelte Polen.

Vi skulle i det Følgende betragte hvert af de tre
Villiestrin for sig, og begynde da med \emph{Driftvillien}.

Mennesket har Drifter ligesom Dyret, men derfor er
Mennesket ikke et Dyr, og de menneskelige Drifter ikke
dyriske. Man bør ikke ringeagte Drifterne, de høre nødvendig
med til vor Menneskenatur; kun naar en menneskelig
Drift begynder at blive dyrisk, har Mennesket Grund
til at blive bange for sig selv. Særkjendet paa den menneskelige
Drift er, at det Sandselige kan optage det Ideelle,
gaae ind under det Fornuftige og suge Kræfter af det Aandelige.
Driften kan omfatte det Laveste og det Høieste; der
gives en Drift til sandselig Nydelse, men der gives ogsaa
en Drift til Guds Rige.

Driften er Villiens frugtbare Skjød, det Hylster og
Svøb, hvoraf den vikler sig ud. Derfor virker Driftvillien
med en saa forunderlig primitiv Naturkraft. Villien kan
ikke gribe uden at gribes af Noget, ikke bevæge sig med
% --- p. 17 ---
\setcounter{footnote}{0}%
\opage{17}Kraft uden inderlig og stærkt at være bevæget af
Noget. Det som fra Begyndelsen af driver og bevæger,
fordi dets Gjenstand fylder Villien, den naturlige Villens
indre Drivende, er Driften.

Driftvillien kan tilegne sig det Fornuftige og fyldes
af Ideer, men derfor bliver den ikke Fornuftvillie; den
bliver kun en \emph{fornuftig} Driftvillie. Den kan have
Salighedens Længsler, tilegne sig Frihedens Forjættelser, ja
selv den dybe Kjærligheds aandige Væsen, og dog bliver
den aldrig Aandsvillie; den er og bliver en aandelig Driftvillie.
Man tænke paa Forholdet mellem Mand og Kvinde,
mellem Moder og Barn, og man vil erkjende, at det Sandselige
kan forklares i det Ideelle, det Aandige, paa Driftlivets
Grundlag. Lidenskaberne kunne være paa een Gang
baade sandselige og ideelle. Digterne prise Kjærligheden,
og vi ere dem taknemmelige. Men Kjærlighed er jo en
Villie, og dog kalder man den ogsaa en Lidenskab. Rigtigt!
Hos Eros er Villien Lidenskab, og Lidenskaben Villie;
saaledes maae det være: det er Driftvillie. De der paastaae,
at denne Driftvillie i sin tilslørede Grund er dyrisk, bagvaske
Eros og haane det Menneskelige. Julie har Villie
til Romeo, hendes Villie er Lidenskab; men denne Lidenskab
er som en aandig Livsluft det høiere Element, hvori hun
voxer; dens Blus og Flamme er Lyset, hvori hun staaer
som forklaret.

Hvor stort Driftvilliens Omfang er, kan sees, naar den
anvendes paa Bestemmelsen af \emph{Folk} og \emph{Mængde}. Hvor
er Folket? I Mængden. Hvor er saa Mængden? 1 % PRINTED AS IS: the figure 1 for I
Folket.
Og dog er Mængden endnu ikke Folket, og Folket mere
end Mængden, skjøndt begge Bestemmelser falde indenfor
Driftvillien. Vi have et storstilet Exempel i det gamle
Testamente. Da Moses var hos Jehova paa Bjerget, slog
Folket sig løs, blev en Hob, blev Mængde. Denne i Lidenskab
bevægede Mængde fordrede af Aaron, at han skulde
gjøre dem Guder, som kunde føre dem an og »gaae for
deres Aasyn.» Det var Driftvillie. Guder af Guld, Guldguder,
det var noget anderledes Virkeligt og Nærværende
% --- p. 18 ---
\setcounter{footnote}{0}%
\opage{18}end hiin Usynlige deroppe i Skyerne. Mængden fordrede
det; den var høirøstet, thi den havde en Villie. Aaron
fandt, at det var en kraftig Villie; saa forvexlede han
Folket med Mængden og støbte Guldkalven. Da Moses
kom, og Aaron saae hans Aasyn, blev han angst og bad:
»ikke blusse min Herres Vrede op; du kjender Folket, at
det hælder til det Onde.» Da viste det sig, hvad Aandsvillie
er. Moses stillede sig for Leirens Port og bød: »hvo
som hører Jehova til, han komme til mig!» Moses satte
Livet ind, thi Hoben var rasende. Men da Kommandostemmen
lød, den myndig bydende, fik Hoben andre Tanker,
en anden Villie: Mængden blev igjen et Folk. Fik Folket
saa med Eet en Aandsvillie? Nei! kun Moses havde Aandsvillie;
men Folket fik ved ham en Villie, som af uvilkaarlig
Drift maatte bøie sig for Aand; det var en aandelig Driftvillie,
og høiere Villie faaer Folket aldrig som Folk.

Gaae vi ind paa de særskilte Former, saa er enhver
aandig Stræben, der oprindelig beroer paa særlig Begavelse,
Talent, Geni, at henføre under Driftvillien; Genivillie er
Driftvillie.

Det er et dybt og sandt og rørende Træk i Øhlenschlægers
Correggio, at den stakkels Antonio, hvis Mod
Mikkel Angelos haarde, uretfærdige Dom har knuust, gaaer
i Rette med sig selv og trøster sig med sin Konstnersamvittighed,
idet han siger: »Gud veed, jeg gjorde det ei
for at glimre. Jeg malede, som Bien danner Celler, som
Fuglen bygger i et Træ sin Rede,» altsaa af indre Trang,
af Konstnerdrift; det er Driftvillie.

Gjælder det om praktisk udadgaaende Virksomhed, om
historiske Bedrifter, hvilken Betydning have da ikke Drifttilskyndelser
og Lidenskaber! Vi see det paa Heltene, en
Kreds af store Skikkelser, som alle staae under Driftvillien.
Hvormeget Fornuftigt, hvormange Ideer og Tanker har
Villien ikke her formaaet at tilegne sig; og dog er det,
naar vi see nærmere til, ikke Fornuftvillie, endnu mindre
Aandsvillie, men en af Ideer beaandet, af Lidenskaber gjennemglødet
Driftvillie, vi have for os. Det er en aandig
% --- p. 19 ---
\setcounter{footnote}{0}%
\opage{19}Naturkraft, en ubetvingelig Attraa, der driver Helten frem;
han vil, fordi han maa ville. Skule Jarl i »Kongsæmnerne»
kan ikke opgive sin Kamp mod Hakon, ikke fraskrive sig
Thronen og gjøre Forlig; og hvorfor ikke? Det er hans
Væsens Trang imod; de store Kongetanker sige nei.

Det er et eget, opløftende Syn at betragte de mange
høist forskjellige Phænomener, som Driftvillien i hver Tidsalder,
hos hvert Kjøn og i enhver Retning kan afsætte;
de ere som Meteorer, de glimre; men jo mere de glimre,
desto lettere ville de briste. Driftvilliens Individer ere med
al deres Villiekraft ikke frie, men naturbundne: de ere i dunkle
Magters Vold. Vel fare de ikke frem i Blinde, thi hvis de
blindthen følge Lidenskabernes Indskydelser, blive excentriske,
altfor kometagtige, saa tabe vi Interessen for dem. Men
selv, naar de ville holde sig i Ideens Spor inden visse
Grændser, er deres Vei den farligste af alle. De have
Reisning, Opsving, Bevægelse og Fart; men Grundlaget
gynger, og Bevægelsen er som en Dands over Afgrunden.
De kunne faae Magt over Andre, men Ingen iblandt dem
alle kan i strengere Forstand faae sig selv i sin Magt.

Det ligger i Driftvillien, at den bevæges af Noget,
hvorover den ikke selv raader.

Som nu i det Erotiske. Forelskelse, siger man, er
Noget, der kommer paa Folk, uden at de egenlig vide,
hvorfra og hvorfor. Det er ikke blot Ungdommen der er
udsat; ingen Alder er ganske sikker: halvgamle Mænd især,
naar det kommer paa dem, blive i Reglen stærkt medtagne.
Men hvad der saadan kommer af sig selv, kan jo gaae
over igjen af sig selv; det er Driftvilliens Luner, Ingen
kan gjøre for det. Især viser sig det Usikkre, naar to
modstridende Følelser eller to Lidenskaber taarne sig mod hinanden,
ligesom to Bølger, og der opstaaer en Brænding.
Den Paagjældende er som i Havsnød; han raader jo ikke
for, hvilken Bølge der skal blive den høieste eller med andre
Ord, hvilken af de to Villier der i ham skal blive den
stærkeste. Sættes den erotiske Villie f. Ex. mod Konstnervillien,
saa har man en skjøn Fortælling om en ridderlig
% --- p. 20 ---
\setcounter{footnote}{0}%
\opage{20}Elsker, som brændte sin Harpe for at frelse sin Elskede fra
at døe i Kulden. Men Forholdet kan let vendes om; thi
mangen poetisk Elsker har ladet sin Elskede fryse ihjel, og
conserveret Harpen. Ingen er ansvarlig, den Troløse saa
lidt som den Trofaste; Ingen kan kaldes til Regnskab, da
jo Alt kommer an paa den stærkere Bølge. Der er en
Romance af Schak Staffeld, »Tilbagekomst»; den handler
om en maurisk Kriger, som vender hjem fra Tagos Bred
med »Kampens Støv i sine Lokker og dens Blod ved sine
Fjed;» han kalder paa sin Elskede: Selindaja! Men hun
gjør Omsvøb; han maa nævne, hvad han vil, Grotten,
Mandeltræet, Gitteret; der er bestandig Noget i Veien, Alt
er blevet anderledes. Tilsidst udbryder han i ædel Harme:
»Ha, Troløse, mens jeg sanked disse femten Hædersar, uden
Blu Du til en Anden Tro og Hjerte givet har;» men hun
svarer ham koldblodig som den, der har en god Samvittighed:
»Hjertet har jeg ei bortgivet, Allah høre mine Ord! Almoralif
har det taget, og hvad kunde jeg derfor?»

At Heltene med al deres Villiesstyrke i Grunden ere
lige saa usikkre paa sig selv som Elskerne, --- forudsat, at
deres Villie kun er en Driftvillie --- kan let bevises. Undersøger
man, hvad det egenlig er, som i en Heltecharakteer
med saa mange Planer, under saa mange Svingninger og
Krydsninger af ydre Omstændigheder, holder det Hele sammen
og grundfæster Eenheden, da vil man ved nærmere Eftersyn
finde, at det ingenlunde er en selvmægtig Villie, men en
eller anden gjennemgaaende, herskende Lidenskab. En saadan
Lidenskab kommer ikke altid tilsyne paa den brusende
Overflade, den holder sig helst under Bølgerne og arbeider
stille i Dybet. Men at det ikke er Helten selv, men Lidenskaben,
der er den egenlige Herre, bliver iøinefaldende,
naar Omstændighederne føie det saa, at der i Sindet vækkes
en ny og mægtig Lidenskab, saa at Dybet røres og de
to Modstandere komme i Strid med hinanden. I et saadant
Tilfælde kan den Stærke pludselig blive svag, faae
Anfald af Vanvid, og Charakteren ligesom gaae i Stykker.
Fra den gamle Historie blot et eneste Exempel: Marcus
% --- p. 21 ---
\setcounter{footnote}{0}%
\opage{21}Antonius. I Triumviratet mellem Octavian, Antonius og
Lepidus, --- til Hævn over Cæsars Mordere, kunde Antonius,
hvad Begavelse og Feltherretalent angik, vistnok i alle
Maader maale sig med Octavian. Ikke desto mindre bukkede
han under, da det siden kom til Fjendskab imellem dem.
Hvorfor? Den almindelige Forklaring er, at Antonius gav
sig til Priis for udsvævende Lidenskaber, medens Octavian
under Kampens Forvirring forstod den Konst at beherske
sig, beregne hvert Skridt paa Sømmelighedens Vei, og holde
Togaens Folder sluttede. Man vilde imidlertid tage meget
feil, om man antog, at Octavian var en Mand, der kun
fulgte Fornuftens Bud, medens Antonius lod sig lede af
blinde Indskydelser. Sandheden er, at begge staae under
Driftvillien. Charakterernes Forskjel er væsenlig at bestemme
ved Forholdet mellem deres Lidenskaber. Octavian, den
siden saa lykkelige Keiser August, havde den Lykke: kun
at være behersket af een Lidenskab, en Lidenskab, der hos
ham var stærk nok til at tvinge og kue de øvrige. Octavians
Sundhed var, at han kun havde een Syge --- Herskesygen.
Antonius derimod havde det Uheld at blive betagen
af to Lidenskaber, befængt med en dobbelt Syge, Herskesyge
og Elskovssyge. Han kom til Ægypten, saae Cleopatra, og
--- hvad kunde han derfor? --- hun erobrede ham, gjorde
ham til sit Bytte. Saa blev Slaget ved Actium afgjørende.
Antonius flygtede, fordi Cleopatra flygtede --- til Ægypten,
og da hun, den Troløse, holdt sig skjult og lod udsprede
det Rygte, at hun var død, tog Antonius sig selv af Dage.
Helten faldt, et ulykkeligt Offer for en lige saa trofast som
utilregnelig Kjærlighed.

Efter sin Natur er Driftvillien den samme, enten den
yttrer sig i Hverdagslivets smaa Forhold eller den gjør
Udslag i store Verdensbegivenheder; men naar vi ville
gjennemskue dens Væsen, er det bedst, for Tydeligheds
Skyld, at see den i store Træk. Det viser sig da, at
Styrkens Hemmelighed er Svaghed, at Selvstændigheden er
uselvstændig, at den opbrusende Frihed bukker under for
Nødvendigheden. Driftvilliens Helte og Heltinder gaae
% --- p. 22 ---
\setcounter{footnote}{0}%
\opage{22}freidigt ombord paa Begivenhedernes Strøm; det tager sig
prægtig ud at see dem i Opseiling: hvert Fartøi gaaer for
fulde Seil, thi Lidenskaberne fylde Seilene. Men her er
farlige Grunde, Compasset har uberegnelig Misviisning, Driftvillien
er som en beruset Styrmand, og Nødankeret mangler.
Saa lide de Skibbrud, de strande paa Skjæbnens Klipper.
Mod Drifternes, Tilbøielighedernes, Lidenskabernes urolige
Element staae Tilværelsens Skranker, de strenge Verdenslove,
Virkelighedens faste Mure; det er Skjæbnens Magt, ---
saa ender det tragisk. At en Enkelt, lykkelig som Octavian,
undgaaer Skibbrudet, er kun en Skuffelse, et Phænomen,
der kun vedkommer det Ydre; i Sagen selv ere de
alle under Skjæbnen.

Det er Brydningen mellem det Urolige og det Faste,
mellem det Driftagtige og det Fornuftige, hvori Selvet løses
op og gaaer tilgrunde. Naar den ydre Modstand hidser
Lidenskaben, naar den indre Syden og Kogen tager til, og
Brændingen stiger, da har Dybet sine egne Tanker, det
tænker paa Undergang; de store Skibe, de stolte Seilere
slaaes til Vrag, Driftvilliens Kræfter splittes. Driftvillien er
altsaa ikke den frie Villie; skal Villien blive fri, maa den
forvandles, hæves op paa et nyt Trin og undergives Lovene.
Vi skulle i det Følgende see et Forsøg.

% --- p. 23 ---
\setcounter{footnote}{0}%
\lecture{Tredie Forelæsning.}\opage{23}

Det er en Forvandling, der foregaaer med Villien,
idet den hæver sig fra sit første, naturlige Trin op paa
sit næste og høiere, og bliver fra Driftvillie til \emph{Fornuftvillie}.
At denne Overgang er mere end en Forandring,
at den, væsenlig seet, beroer paa en Forvandling, maa
udtrykkelig fremhæves.

Driftvillien kan, som i vort Foregaaende er viist,
klares, ændres, idealiseres, kan fra at være en sandselig
blind Villie blive til en seende Villie, en Villie med Tanker,
Ideer, Planer, det er en Fremgang til det Bedre, en Forædling,
en stor Forandring, men det er endnu ingen Forvandling.
Det inderste Bevægende er Drift og bliver
Drift, det Hele falder indenfor ufrie, dunkle Naturtilskyndelsers
Omraade. Saalænge Driften raader, er Lidenskaben
den egenlige Kraftmaaler; der gives ikke blot opblussende,
hidsige, hæftig udfarende Lidenskaber, der gives ogsaa
sindige, tause Lidenskaber, rugende, drømmende, ulmende
Lidenskaber. Men indenfor den hele Kreds af umiddelbare
Naturtilskyndelser, de være lavere eller høiere, hidsige
eller sindige, gjælder det: jo større Lidenskab, desto stærkere
Villie, Lidenskabens Styrke er Betingelsen for Villiens
Kraft.

Først idet Villien gaaer over fra at være fornuftig
Driftvillie til at blive, hvad vi i egenlig For%
% --- p. 24 ---
\setcounter{footnote}{0}%
\opage{24}stand maa kalde Fornuftvillie, skeer Forvandlingen. Alt
bliver fra Grunden af anderledes. Drifterne udryddes ikke,
men Magten tages fra dem, de tæmmes; de naturlige
Tilskyndelser træde i Fornuftens Tjeneste, og maae forpligte
sig til ubetinget Lydighed. Nye Bevæggrunde, en
heel Besætning af nye Motiver, komme til Roret; det er,
bogstavelig talt, et nyt Dynasti, der bestiger Thronen.
Driftens Sted træder Indsigten; ikke dunkle Rørelser, men
klare Tanker, ikke gjøglende Haabs tillokkende Billeder,
men Erkjendelse af det Almeengyldige, fast Overbeviisning
om det i og for sig Sande og Rette er det, som fra nu
af skal bevæge Villien. Her spørges ikke om Lyst og
Ulyst, ikke om Tilbøieligheder, men om Pligt og Ret, thi
Pligt og Ret ere personlige Udtryk for uafviselige Love.
Her er i denne Anskuelse en Høihed, som imponerer; en
kjølig Helligdom, en Kuppel med evige Stjerner hvælver
sig over det Timelige. Mennesket fatter sin væsenlige
Bestemmelse  og faaer Agtelse for sig selv. Men i det
Tiltrækkende er der tillige noget Frastødende, Noget, som
kan opvække Gysen. For at gjøre Villien fri, maa Fornuften
begynde med at overskjære de Baand, som binde den til
Drifternes moderlige Rod; hvert Nerverør, hver Trevl og
Traad fra Moderstammen maa sønderrives. Skal Fornuften
raade, saa maa den raade i Et og Alt. En saadan Selvhersker
gaaer ikke paa Accord med Hjertets Rørelser, vil
ikke parlamentere med lunefulde Tilbøieligheder, ikke erkjende
noget Følelsens Veto: Fornuften er en Souverain
med absolut Monarchie.

Det vil heraf forstaaes, at det ikke er Handlingens
Udvortes, hvorpaa det væsenlig kommer an, men dens
Indvortes. En Handling, der udspringer af en fornuftig
Driftvillie, kan i sit Ydre have en paafaldende Lighed
med en Handling, der er udsprunget af reen Fornuftvillie:
Forskjellen ligger i det indre Bevægende. Det er, naar
der tales om historiske Personer, ofte vanskeligt at afsige
en sikker Dom. En Agamemnon offrer sin Datter, en
Junius Brutus »slagter sine Sønner paa Borgerfrihedens
% --- p. 25 ---
\setcounter{footnote}{0}%
\opage{25}Alter»: hvad er det nu egenlig for Handlinger? De, som
strax ere ved Haanden med deres: »blinde Fanatikere,
Umennesker, unaturlige Fædre!» forraade just ikke nogen
dyb Indsigt enten i Driftvilliens eller i Fornuftvilliens
Væsen; thi det er vist, at der gives Optrin, Collisioner,
forfærdelige Sammenstød af modstridende Fordringer, hvorunder
den kolde Fornuft kan tyrannisere det naturlige
Hjertes Følelser fuldt saa stærkt som en i Lidenskab
glødende Driftvillie.

De anførte Handlinger kunne, naar de sees paa Afstand,
og betragtes i og for sig, være Fornufthandlinger,
thi det er jo Fornuftens Bud, at alt det Individuelle skal
offres for det Almene. Ja, de \emph{kunne} være det. Men
see vi nærmere til, saa staaer hverken Agamemnon eller
Brutus paa den rene Fornufts Standpunkt; de ere begge
optagne af endelige, om end paa Samfundets Almene
sigtende, Formaal: deres Villie er en potenseret Driftvillie.
Paa de endelige Formaals Omraade skal Fornuften vel
regulere Driften, men for Resten skaane den, saa at den
inden visse Grændser faaer Lov til at være sig selv i al
Naturlighed. Fra dette Synspunkt falder der en mørk
Skygge over de store Skikkelser, Agamemnon og Brutus;
thi Faderhjertet har lige saa fuldt sin ubestridelige Ret,
som Samfundets Almene har Krav paa Kongen og Statsmanden.
I Driftvilliens Navn fordre vi, og med Føie,
at Fornuften skal være menneskelig, og dens strenge Bud
lade sig forlige med Humanitet. Det fordrede Grækerne
--- dog Spartanerne fraregnede --- ogsaa; men ikke Romerne.


Den mærkelige Forskjel, der, i det Hele taget, er
mellem det græske og det romerske Sind, mellem Grækernes
livsvarme, geniale Driftvillie, og den kolde, men følgestrenge
Romervillie, er kjendelig selv af de tilfældig valgte
Exempler. Det Haarde, Oprørende i Agamemnons Offerhandling
har det græske Sagn paa eiendommelig Maade
vidst at formilde og forsone med skjøn Humanitet. Artemis,
den fortørnede Gudinde, kommer Menneskefølelsen
% --- p. 26 ---
\setcounter{footnote}{0}%
\opage{26}til Hjælp, hun fordrer sit Offer, men hun frelser det ogsaa;
hun frelser Iphigenia, fører hende bort til Tauris og
kaarer hende til sin Præstinde. Indtrykket af den guddommelige
Mildhed borttager det Oprørende og virker som
et Idealitetens Bad, hvori den naturlige Følelse luttres.
Anderledes med det romerske Optrin. Vistnok havde
Brutus's Sønner forbrudt sig imod Republikken, forsaavidt
de hørte til de unge Sammensvorne, der vilde gjenindføre
Kongedømmet, og Brutus, »Fædrelandsbefrieren», kunde ikke
uden Selvmodsigelse stemme for deres Benaadning. Men
Maaden, hvorpaa Henrettelsen bliver skildret, den forstenende
Kulde, hvormed Faderen, uden engang at vende
sit Ansigt bort, seer Sønnernes Hoved falde, er oprørende
for al Humanitet, et Symbol paa Romervilliens Foragt for
naturlig, menneskelig Følelse.

Denne hjerteløse Haardhed, denne ubøielige, al naturlig
Følelse kvælende Følgestrenghed kan let imponere Betragteren,
forvirre hans Dom og forlede ham til at ansee
Romervillien ikke blot for mere fornuftig end Grækervillien,
men endog for en reen og stor og ophøiet Fornuftvillie,
en mønsterværdig, systematisk gjennemført Statsmandsvillie.
Det er en Indbildning. Ganske vist er Romervillien
anderledes systematisk conseqvent end Grækervillien,
og derfor ogsaa stærkere. Den græske Villie i sin foraarsfriske
Ungdommelighed har Kræfter, store, geniale
Naturkræfter; men Romervillien er mere end en Udstraaling
af Kræfter, den er en egen Kræfternes Eenhed, en
Magt, en verdenserobrende, folkeopslugende Magtvillie. At
denne Magtvillie, langt fra at være, hvad vi i egenlig Forstand
maa kalde Fornuftvillie, men en aandløs-aandig Naturvillie,
en ulvegraadig, umættelig Driftvillie, som, naar den
har fortæret Alt, maa ende med at fortære sig selv, trænger
i denne Sammenhæng ikke til Beviis: Historien fører
Beviset.

Det moderne Begreb om Fornuft og Fornuftvillie er
endnu ikke klarere, i praktisk Forstand neppe saa klart,
som det var hos Oldtidens Vise. Vi fordre Fornuft og
% --- p. 27 ---
\setcounter{footnote}{0}%
\opage{27}vi fordre Humanitet; men Spørgsmaalet er, om disse to
Fordringer virkelig lade sig forene, saa sandt Fornuften,
vel at mærke, skal være vort Et og Alt, den eneste,
høieste styrende Magt i vor humane Verden.

Vi fordre, at Fornuften skal beherske os og dog
skaane vort Hjertes Tilbøieligheder; vi mene, at det vel
lader sig gjøre, da vore ædle Tilbøieligheder jo i Grunden
vidne om fornuftige Anlæg og altsaa maae være villige til
at gaae ind under Fornuftvilliens ubetingede Herredømme:
det turde dog være en Feiltagelse. Vi fordre af Fornuften,
at den ikke blot skal være streng, men ogsaa billig, skal
tage Hensyn til vore Skrøbeligheder, kort: opdrage os
som en streng, men kjærlig Fader opdrager sine Børn.
Men Lignelsen halter. Fornuften er ingen kjærlig Fader,
men en ubøielig Hersker, kold, uforanderlig som de evige Love. Naar Faderen revser sin lille Dreng, saa gjør den
Lille Omvendelse, lover Bod og Bedring: »vil være artig,
vil være sød!» Saa gaaer Uveiret over og det kommer
til en Forsoning, skjøndt Faderen veed, at den angerfulde
Synder lige saa let kan glemme Løftet, som han glemmer
det snertende Riis. Men saaledes kan Fornuften ikke
bebandle % PRINTED AS IS (behandle)
os Myndige, kan umulig paa een Gang baade
rise og tilgive os Fornuftige. Vi fordre, at Fornuften
skal respectere vor Individualitet, skal vise Imødekommen,
være en føielig, human og medgjørlig Fornuft: paa disse
constitutionelle Betingelser skal den faae Lov at regjere.
Men paa disse Betingelser vil Fornuften ikke indlade sig.
Og hvorfor ikke? Fordi den ikke kan det uden at modsige
sig selv og blive ufornuftig.

Er Fornuften da despotisk, inhuman og grusom?
Ingenlunde. Det være langt fra os at ville nedsætte Fornuften,
tale ilde om den og forsøge at sætte en Plet paa
dens rene, ophøiede Charakteer. Vi ville tvertimod lade
Fornuften selv komme til Orde. Den siger til Mennesket:
Du har at vælge. Vil Du blive staaende paa Driftvilliens
Standpunkt, velan! Jeg skal tjene Dig med Alt, hvad
jeg kan, med Raad og Formaning, Underviisning og Op%
% --- p. 28 ---
\setcounter{footnote}{0}%
\opage{28}lysning i enhver Henseende. Men saa maa Du vide, at
jeg kun er Din Tjener, ikke Din Herre. Du kan følge
mine Raad, om Du saa behager, og Du kan forkaste dem
om Du vil: jeg indestaaer ikke for Følgerne. Blive dine
Følelser splidagtige, gjøre dine Lidenskaber Oprør, kommer
dit Sind i Kog og der dannes en Malstrøm, som suger
Dig til sig og maler Dig ned i en Afgrund, det bliver
Din egen Sag: jeg er uden Ansvar. Vil Du derimod udfries
af din naturbundne Tilstand og løftes op paa Selvstændighedens
Trin, saa maae vi vende Forholdet om;
saa bliver det mig, der skal byde, og Dig, der skal lyde,
saa er jeg ikke din Tjener, men ubetinget din Herre:
vælg nu! Det kan man da ikke kalde inhumant; det er,
som mig synes, ret fornuftig talt af Fornuften.

Det Eneste, man kan befrygte, er at Fornuften, naar
den bliver saa enevældig, vil glemme at benytte sig af
den skjønne kongelige Ret, man kalder Benaadningsret.
Og saaledes er det. Fornuften er blandt alle absolute
Herskere den Eneste, som ikke kjender til Benaadningsret.
Fornuften er det selvløse Almene, hvis Væsen og Indhold
gaaer op i Lovenes Sammenhæng. Fornuftvillien er Selvets
egen Villie, forsaavidt det tænkende Selv alene vil, at Nødvendighedslovene
opfyldes. Frigjørelsen bestaaer i, at Selvet
vil, hvad det af Overbeviisning maa ville: opløse sig selv
i det Almene og Selvløse.

Paa Driftvilliens Trin befinder Individet sig vel, naar
dets individuelle Krav, de være materielle eller ideelle, tilfredsstilles;
derfor er Sundhed et Gode, Sygdom et Onde,
Rigdom et Gode, Fattigdom et Onde o. s v. Paa den % PRINTED AS IS: no stop after s
Art Tilfredsstillelse kan Fornuftvillien ikke indlade sig.
Alle baade indvortes og udvortes, baade usandselige og
sandselige Goder, forsaavidt de blot tjene til at mætte
den umiddelbare Individualitet, ere af underordnet Betydning.
Man behøver just ikke at foragte dem; men
de underligge Tilfældighed, betinges af Lykke, og maae
desaarsag henregnes til det i sig selv Ligegyldige. Den
Vise er hævet over Lykkens Omskiftelser; han haaber Intet,
% --- p. 29 ---
\setcounter{footnote}{0}%
\opage{29}thi han frygter Intet: det er det gamle stoiske Ideal, og
dette Fornuftideal er endnu det eneste conseqvente. Den
Vise er sig selv nok; han vil, hvad Skjæbnen vil, hans
Liv er en Skjæbnens Aabenbarelse: saa er han stærk som
Skjæbnen, klippefast, urokkelig som Skjæbnen.

Betragte vi nu dette Fornuftideal --- og der gives
intet andet --- noget nøiere, saa viser det sig, at den Charakterens
lovpriste Overeensstemmelse med sig selv, hvoraf
der i Oldtid og Nutid er gjort saameget Væsen, opløser
sig i Modsigelse paa Modsigelse, og at disse Modsigelser
ingenlunde ligge i den evige Fornuft, men i Menneskenes
egen Misforstaaelse af, hvad Fornuften er, og hvortil den
skal bruges.

Det er da for det Første let at indsee, at det opstillede
Ideal ikke kan virkeliggjøres. De høittravende
Talemaader, at den Vise er sig selv nok, er lykkelig selv
i en gloende Ovn, er Konge, Præst, Spaamand og Digter,
maae tages for hvad de ere --- Talemaader.
Siger man:
et saadant Ideal stilles op som et Maal, hvortil der i det
Uendelige stræbes; det naaes kun i uendelig Tilnærmelse:
da er det ikke vanskeligt at overbevise sig om, at Tilnærmelsen
indeholder en Modsigelse. Saalænge nemlig
alt det, den naturlige Individualitet maa betragte som
væsenlige Goder: Sundhed, Begavelse, Rigdom, Anseelse,
Magt, Indflydelse o. s. v. endnu ikke er nedsat til at være
noget aldeles Ligegyldigt, altsaa saalænge det indre Bevægende
endnu vedbliver at være en umiddelbar Drift,
saalænge vedbliver Mennesket at staae under Driftvillien.
Der kan paa dette Trin saa lidt være Tale om en Tilnærmelse til det rene Fornuftideal, som der ved at reise
Jorden rundt kan være Tale om en Tilnærmelse til Maanen.
Gjøres der Alvor af, at den bevidste Fornuft skal ved alle
Handlinger være det indre Bevægende, saa at Ingen tør
attraae Noget, med mindre Fornuften byder det, Ingen
hungre eller tørste, elske eller hade uden ifølge Fornuftens
udtrykkelige Befaling, saa begynder det Umenneskelige.
% --- p. 30 ---
\setcounter{footnote}{0}%
\opage{30}Men et i sit inderste Væsen umenneskeligt Maal kan ikke
være et Ideal for Mennesker.

Modsigelsen kommer tilsyne i alle virkelige Handlinger.
At de timelige Goder uden Undtagelse ere ligegyldige, vil
med andre Ord sige, at alle endelige Formaal ere ligegyldige.
Den Vise maa kunne finde sig tilrette i alle
Stillinger, være virksom under alle Forhold. Fornuftvillien
maa virke, dens kraftige Væsen er Virken. Den Vise kan
f. Ex. være Feltherre. Saa leverer han da et Slag.
Hvorfor? For at vinde det! Ja, det mener en Daare,
et Individ, der endnu befinder sig under Driftvillien. Nei,
den Vise leverer Slaget for Fornuftens Skyld, ene og alene
for Fornuftens Skyld. Om han vinder eller taber, kommer
derfra med Ære eller med Skam, er ham ligegyldigt,
ganske og aldeles ligegyldigt. Men naar alle udvortes
Øiemed ere ligegyldige, saa maae de tilsvarende Handlinger
jo ogsaa være ligegyldige. Hvorfor vil da den Vise nedlade
sig til at handle?

Den Vise virker og handler, ikke for andre Menneskers
Skyld, ikke for Folkets, ikke for Statens Skyld, men alene
for sin egen Skyld, at han i sin Væren og Virken kan
vide sig i Overeensstemmelse med sig selv og Fornuften.
Heraf følger, at hans Maalestok for Dyd og Last maa
blive en anden end den, der justeres for det borgerlige
Samfund. Da nemlig enhver udvortes Handling i sig selv
er ligegyldig, maa Handlingens Beskaffenhed, dens Værd
eller Uværd, udelukkende beroe paa Charakterens Overeensstemmelse
med sig selv hos den Handlende. Den
Vise kan i Medfør af sin Selvovereensstemmelsesfordring
være beføjet til at begaae Handlinger, der i Menneskenes
Øjne see ud som Forbrydelser; han kan røve, myrde,
stjæle, lyve, Altsammen i Conseqvents af den rene Fornuft.
I slige Henseender er Intet forbudt, undtagen det Umulige.

Men, indvender man, naar det rene Fornuftmenneske
kan blive et saadant Uhyre, hvorledes kunne da Charakterer
saa ædle og ophøiede som en Sokrates, en Epiktet,
en Marcus Aurelius, regnes iblandt Fornuftmennesker?
% --- p. 31 ---
\setcounter{footnote}{0}%
\opage{31}Vi svare: De staae paa Fornuftvilliens Trin, forsaavidt det
indre Bestemmende ikke er en umiddelbar og dunkel Drift,
men en i Tanken klaret Sandhed. Men ved en ethisk
Selvskuffelse have de gjort Fornuften til noget Mere end
den er, til et aandigt Selvvæsen, en villende Guddom; de
have af deres eget rige Indre lagt Noget ind i Fornuften,
som umuligt kan udledes af dens Væsen. Det er Charakterer
med en dæmrende, religiøs Baggrund: den personlige
Forudsætning er: at Selvet er evigt. En saadan
Forudsætning er vel ikke i Strid med Fornuften, men
--- det maa atter og atter indskærpes --- den kan ikke begrundes
af Fornuften, ikke afledes af det rene Fornuftbegreb.
Fornuftbegrebet, fastholdt for sig, er Begrebet
af det selvløse Almene: Fornuftgrunde ere Nødvendighedsgrunde,
ikke Frihedsgrunde.

Den Frihedsbestemmelse, som kan udledes af den
rene Fornuft, er saa selvmodsigende som muligt. Den
Vise har en Villie, thi han har et Selv. Saa begynder
han da med at luttre dette Selv, gyde det i Fornuftens
Smeltedigel for at befrie det fra Slagger og urene Indblandinger.
Han samler sig og tager fat med en Villieskraft,
en Energie, som var dette Selv bestemt for en heel
Evighed. Men netop her naaer Modsigelsen sit Høieste;
den culminerer, saa at sige, i Fornuftens Zenith. Thi
Selvet, seet i dette Lys, er ikke evigt; det er forgængeligt
som Tingene i Naturen. Det existerende Selv er individuelt,
og det Individuelle løses op og forgaaer: kun
Fornuften er evig, evig som Lovene. Naar den Vise har
spillet sin Rolle paa Verdenstheatret, ligegyldigt, om det
er Tiggerens, Skrædderens eller Keiserens Rolle, og seer, at
Rollen er ude, saa træder han af. Han behøver ikke at
vente paa naturlig Afløsning, han kan afløse sig selv, han
kan med Frihed slukke sit Lys og gaae bort: at give sig
selv Døden er Udtrykket for den høieste Skjæbnevillie, den
Villiesbestemmelse, hvori Selvet selv løser sig op, den
formelle Friheds sidste Fornufthandling.

Men naar Menneskets Selv er forgængeligt som Alt i
% --- p. 32 ---
\setcounter{footnote}{0}%
\opage{32}Verden, og det Forgængelige er et Ligegyldigt, saa er det
dog vel heraf indlysende, at ogsaa Selvet er et Ligegyldigt,
dets Selvstændighed og Overeensstemmelse med sig selv
i Grunden et Ligegyldigt: og dog var det jo netop for
Selvets og Selvstændighedens Skyld --- altsaa for et Ligegyldigt!
--- at der i sidste Instants skulde leves, virkes
og handles.
Dette er den rene Fornuftvillies i Et og Alt
gjennemførte Conseqvents.

Den moderne Bevidsthed indrømmer, at de ægte
Stoikere have indviklet sig i Modsigelse, men maa, naturligviis!
give dem Ret i at det ethiske Fornuftideal er det
eneste Sande. Sagen er, at hine gamle Tænkere vare
eensidige. Deres Verdensanskuelse var kold og gold: Tilværelsens
Indhold, Livstræets frodige Væxt, dets grenede
Stamme, det legende Lys, Individualitetens yppige Skud
med Løv og Blomster i Duft og Farve, alt dette savnes;
alt dette blev overseet, thi de gamle Tænkere vare, som
sagt, eensidige. Men de nye Tænkere ere alsidige, de
tage Alting med, det Mindste med det Største, og oversee
Intet. Hvorledes bære de sig ad? Det synes indviklet
og gaaer dog ligefrem: man gjøre frisk væk Natur til
Fornuft, og Livsindholdet --- det Concrete --- kommer af
sig selv. Fornuften er naturlig i Følelser, Drifter og Lidenskaber;
Naturen er fornuftig i Tanker, Ideer og Love.
Fornuftens Krav er Naturens Ret: Fornuften i Kjød og
Blod, Naturens Gjenindsættelse i sin Ret, er det sande
Fremtidsideal, den dybere Mening i »Kjødets Evangelium».

Saavidt kunde de Gamle ikke naae, dertil var deres
Enfoldighed endnu for stor. Vel begyndte ogsaa de med
en almindelig Eenhed af Fornuft og Natur: at leve overeensstemmende
med Fornuften var hos dem at leve over-
eensstemmende med Naturen. Men dette gjaldt kun i al
Almindelighed; saasnart de kom ind paa det Enkelte, indtraadte
Spliden. At gjøre det Timelige evigt og det Evige
Timeligt, det Ligegyldige Væsenligt og det Væsenlige ligegyldigt,
er en Opgave, hvis Løsning overgik de gamle
% --- p. 33 ---
\setcounter{footnote}{0}%
\opage{33}Tænkeres Kræfter, og maa desaarsag være forbeholdt de
nye Tænkere.

Fornuftens Herredømme og Kjødets Evangelium!
Tør vi nu ogsaa give os hen til det lyse Haab, at en saa
kolossal Modsigelse virkelig vil blive løst? Istedenfor at
gaae fremad, gaae vi jo tilbage; istedenfor at hæve os
op paa Fornuftvilliens høiere Trin, synke vi umærkelig
ned under Driftvillien. Fornuften er ikke moderne, saa
lidt som Naturen; det er slet ikke sagt, at den moderne
Bevidsthed og den evige Fornuft ere Et og det Samme.

Fornuftvilliens Trin er et nødvendigt Overgangstrin,
og Erkjendelsen af den praktiske Fornuft en væsenlig
Betingelse for en kraftig Villie. Det Rensende og Opdragende
ligger i, at det menneskelige Selv lærer at bøie
sig for noget Almeengyldigt, lærer at resignere og igjennem
uendelig Resignation modtager et Indtryk af noget Evigt.
Men Fornuften kan ikke lære Mennesket at opfatte sit
eget individuelle Selv som evigt. Dersom den kunde det, saa
maatte der i Videnskaben gives et fuldgyldigt, uigjendriveligt
Beviis for Sjælens Udødelighed. Men at et saadant Beviis
hverken er givet eller nogensinde vil kunne gives, er en
klar Sag, naar vi først tilfulde have lært at indsee, hvad
Fornuft er, og hvorfor dens Evighed maa være selvløs:
Udødelighedens Mulighed kan bevises, men dens Virkelighed
kan ikke bevises.

Vil man bevise Sjælens Udødelighed af moralske
Grunde og beraabe sig paa, at en fuldstændig Udvikling
af vor frie sædelige Personlighed vilde være umulig, dersom
vor Sjæl ikke var udødelig: da overskrider man Fornuftens
Grændser, overanstrænger dens Kræfter og gjør en
personlig Indsats, for hvis almene Gyldighed den rene
Fornuft hverken kan eller vil gaae i Borgen. Saa bliver
det ved det Gamle: Fornuften, det almene Væsen med
Ideer og Love, er evig; Menneskets Individualitet med sin
organisk betingede Selvbevidsthed, sine Drifter og Livskrav
er timelig. Lægges den afgjørende Vægt paa det Sidste,
% --- p. 34 ---
\setcounter{footnote}{0}%
\opage{34}saa hersker Driftvillien; gaaer Selvbestemmelsen ud fra
det Første, saa hersker Fornuftvillien.

Driftvillie er et Trin for sig og Fornuftvillie et Trin
for sig. De to Trin kunne ikke skrues sammen til eet.
Man kan ikke tage Stykker ud af Driftvillien, for at bøde
paa Fornuftvillien, saa lidt som omvendt. Den personlige
Selvstændighed lader sig ikke consolidere ved at lime
noget Individuelt, som ikke er evigt, sammen med noget
Evigt, som ikke er individuelt. Villiens Frihed er intet
Flikværk: man kan flikke paa en gammel Sko, men man
kan ikke sammenflikke Idealer. Skal Fornuftvillien forenes
med Driftvillien, da maae de begge forvandles og optages
i et tredie Høiere: dette tredie Høiere er Aandsvillien.

% --- p. 35 ---
\setcounter{footnote}{0}%
\lecture{Fjerde Forelæsning.}\opage{35}

Vi leve i Opdagelsernes Tid; her spørges og høres
bestandig Nyt: nye Maskiner, nye Stjerner, nye Solsystemer,
nye Svovlstikker! hvo kan opregne det alt; Intet i Stort
eller Smaat unddrager sig Forskerens Blik. Det mylrer
frem, det hidtil Ubekjendte, fra Havets Grunde, fra Himmelrummets
Dybder og fra Jordens Indvolde. Det Spredte
samles, ordnes, sammenlignes, tydes, vurderes, examineres.
Det er Virkelighedens Indhold, Verdens Sammenhæng, Tingenes
Kræfter og Lovenes Rige, som den tænkende, iagttagende
Fornuft utrættelig gjennemforsker.

De store Opdagelsesreiser gaae ikke blot udefter men
ogsaa indefter. Fornuften dvæler ikke ved det blot Sandselige;
den spørger om det Usandselige, den vil i det
Timelige opdage det Evige, i alt det Vordende erkjende det
Værende. Hvad er saa dette evigt Værende? Er det Materiens
Grunddele, uforanderlige Atomer, hvoraf alle Ting
sammensættes for igjen at opløses? Eller er det et overbevidst
og just derfor ubevidst Tankevæsen, Naturens blinde,
iboende Viisdom, som virker og væver i alle Phænomener?
Eller er Hemmeligheden maaskee et uendeligt Intet, det
Intet, som jævner Alt, den mørke Nat, som gjemmer Alt,
den hellige Taushed, hvori alle Røster forstumme, det
overfyldte-tomme Svælg, det evige Ginnungagap, hvoraf Alt
gaaer frem, og hvortil Alt, det Hele, tilsidst skal vende
tilbage? Hvad er Natur, hvad er Aand, hvor er det Baand,
som binder dem sammen? Her gjør Fornuften Opdagelse
paa Opdagelse; det ene Verdenssyn jager det andet, den
% --- p. 36 ---
\setcounter{footnote}{0}%
\opage{36}ene Philosophie fortrænger den anden, den ene Tænker overbyder
den anden. Materialisme, Spiritualisme, Realisme,
Idealisme, Pantheisme, Atheisme, Nihilisme: det er et
broget, vidtløftigt Prøvekaart der foreligger. Og hvad beviser
saa dette? At her er mange, og, naar vi see nøiere
til, sammenhængende Forsøg, der alle stræbe til eet Maal.
Fornuften vil alsidigt klare sit Syn; den frie Tanke vil
ikke lade sig binde til noget Dogme, ikke vige for noget
Magtsprog, ikke bøie af for nogen vilkaarlig opstillet Skranke.
--- Det er Ret, saaledes bør det være. Vi skulle ikke
sætte Fornuftens Lys under Traditionernes Lampeskærm,
ikke lænkebinde den frie Tanke, ikke ønske, at det dog
nu omsider maatte have en Ende med de mange og store
Opdagelser. Tvertimod! Jo flere Fornuftopdagelser der
gjøres, de tilhøre nu den udvortes eller den indvortes
Verden, desto klarere maa jo den vidende Bevidsthed blive.

En Opdagelse staaer imidlertid tilbage, en eneste,
men meget væsenlig, en Opdagelse, som bestandig bliver
overhalet i Lærdom, og, trods den megen Lærdom, flygtig
behandlet; dens Gaade ligger i det Spørgsmaal: hvad er
et \emph{Menneske}?

»Et Menneske,« svarer den moderne Viden, »er et
heldig organiseret Dyr; det er tydeligt af Jordklodens
Udviklingshistorie. Mennesket nedstammer fra en eller
anden Dyreform i Urverdenen, rimeligviis fra en gammel
prægnant, tvetydig Abe.« --- Dersom der ved denne Nedstammen
kun skal tænkes paa et zoologisk-naturligt Overgangsled,
saa vedkommer Undersøgelsen udelukkende Naturlæren,
og vi kunne lade Hypothesen gjælde for, hvad den
er --- en Hypothese. Men, hvis Afstamningen skal forstaaes
saaledes, at Menneskenaturen oprindelig er dyrisk, en
Abenatur, som i tusinde Gange Aartusinders Løb gjennem
utallige Ændringer og gradvise Overgange har arbeidet sig
op til at blive et selvbevidst Fornuftvæsen, et Væsen, hvis
dunkle, naturbundne Abevillie lidt efter lidt er blevet til
en menneskelig klar og fri Fornuftvillie: da bliver Opdagelsen
philosophisk, og gaaer ind under Aandslæren.

% --- p. 37 ---
\setcounter{footnote}{0}%
\opage{37}Det Philosophiske, siger man, er saa vanskeligt at
forstaae; det er kun for de gode Hoveder. Ja, det Philosophiske,
vi her have for os, er ganske vist vanskeligt at
forstaae. En Abenatur, der er sig sin Fornuft bevidst,
videnskabelige Aber, Aber, der philosophere! det Philosophiske
i en saadan Philosophie ligger over almindelig Menneskeforstand;
det kan kun forstaaes --- i klar Logik
--- af Familiens gode Hoveder.

Vilde man fra dette Synspunkt indlade sig paa at
undersøge Værdien af det Udbytte, Mennesket har vundet
igjennem den gradvise Udvikling, da er her ganske vist,
fra een Side seet, forbausende Resultater at opvise; men
spørge vi, hvorvidt alt dette egenlig har været til Gavn for
Slægtens Frihed og Lykke, da opstaae Betænkeligheder,
ikke mindst med Hensyn paa Villien. Dyret vil, hvad det
maa ville, det nøies med sin Driftvillie; Mennesket har
paataget sig at ville, baade hvad Driften vil, og hvad Fornuften
vil: det er en stor Besværlighed. Aben er lykkelig,
den vil altid helst, hvad den bør; Mennesket maa gjøre
en Forskjel imellem, hvad det helst vil, og hvad det bør
ville; det er slet ikke lykkeligt. Ved sin beundringsværdige
Selvudvikling, sin megen Kultur og Civilisation har
Mennesket paadraget sig den store Ulykke at blive bebyrdet
med to Villier, som ikke kunne forliges.

Det kommer der ud af de mange Opdagelser. Mennesket
har opdaget sig selv og i Selvet opdaget Spliden;
Spliden er i dets Indre, Modsigelsen gaaer til Marv og
Been, thi Selvet er dyrisk, og dog fornuftigt. Aben, der
aldrig har indladt sig paa at opdage sig selv, har lykkelig
undgaaet Spliden. Dyret har et Naturselv, men seer det
aldrig, og slipper let fra alle Bryderier med Modsigelsen.
Mennesket, der har dreiet sit Øie indad og seet sit bevidste
Selv, er blevet et Offer for Modsigelsen. Af Modsigelsen
udspringer Tvivlen, og med Tvivlen følger Vildfarelsen:
Mennesket har opdaget Tvivlen og i Tvivlen Vildfarelsens
Mulighed. Der er i ældre Dage holdt udførligt Register
over de mange positive, baade legemlige og sjælelige, For%
% --- p. 38 ---
\setcounter{footnote}{0}%
\opage{38}trin, som skulde udmærke Mennesket frem for Dyret. Det
Positive i alle disse Fortrin har den nyere Kritik saa godt
som udslettet; de negative Fortrin derimod har Kritiken
ikke kunnet udslette: Mennesket kan fordærve sit Liv, forvrænge
sin Opgave, fare vild og forfeile sit Maal; det Fortrin
har Dyret ikke. Mennesket har opdaget Dyden, thi
det har opdaget Lasten, og det opdagede Lasten just den
Dag, da det opdagede Pligten. Dyret er dydigt, men uden
at vide deraf, det handler i Overeensstemmelse med sin
Natur, det kjender ikke Lasten, thi det kjender ikke Pligten.
Mennesket har gjort Bekjendtskab med Pligten, ak!
det gjør et betænkeligt Skaar i Dyden: dersom der ingen
Pligter var, saa vare vi alle dydige. Mennesket har opdaget
Døden, en uhyggelig Opdagelse. At Dyret maa døe
lige saa vel som Mennesket, er vist; men Døden som Død
er ikke til for Dyret. Døden er i sig selv som et Intet;
men det er forbeholdt det selvbevidste Væsen at tænke
dette Intet, møde dette Intet som en Magt og give det
Tilværelse i Tanken. Dyret lever i Øieblikket og nyder i
Øieblikket; hvert givet Øieblik faaer Sit, og det næste
Øieblik er ikke til. Mennesket seer i det nærværende
Øieblik Muligheden af et næste Øieblik, og i denne Mulighed
seer det Døden. Saa kan Synet af et næste Øieblik
pludselig forstyrre Nydelsen, den manende Skygge kan vise
sig midt under Festen, som Bancos Aand ved Macbeths
Gjæstebud. Mennesket har opdaget Evigheden; det er
endnu det Allerværste. Denne Opdagelse sætter Kronen
paa alle dets Fortrin, dog, vel at mærke, negativt, thi det
er en Tornekrone. Ved at blive sig det Evige bevidst opdager
Mennesket, at Selvmodsigelsen i dets Væsen er uløselig:
det Evige forpligter, og Modsigelsen er i Forpligtelsen.


Hvortil er Mennesket forpligtet? Til at gjøre det
Gode --- for det Godes Skyld. Heri ligger, at det Gode
ikke tør maales med Endelighedens Alenmaal og sættes i
Klasse med det blot Nyttige, det i Øieblikket Tjenlige;
en blank og bar Nyttighedslære er i sædelig Forstand en
% --- p. 39 ---
\setcounter{footnote}{0}%
\opage{39}fræk Fornegtelse af det Evige. Men naar det Gode skal
villes, fordi det i sig selv er godt, saa kan et Menneske
jo i timelig Forstand blive ulykkeligt, alene fordi han vil
det Gode. Naar da vor Nedstammen fra Dyret bliver forstaaet
saaledes, at Mennesket døer som et Dyr, og Meningen
er, at det menneskelige Selv lige saa lidt er evigt som
det dyriske: saa maa jo det Menneske, der vil det Gode
for det Godes Skyld, indvikle sig i en uopløselig Selvmodsigelse.


Det Gode, siger man, har sin Løn i sig selv; Mennesket,
der føler Forpligtelsen, faaer en forhøiet Selvfølelse,
en uvilkaarlig Fornemmelse af, at det Evige er et Nærværende.
Det er en tvetydig Tale. Den kan være sand,
saalænge Spørgsmaalet om Selvets Evighed endnu staaer
hen ubesvaret. Det er et daarligt Tegn paa Sands for det
Gode, naar et Menneske vil værge sig mod Forpligtelsens
Krav ved at kræve Beviis for, at en timelig Opoffrelse virkelig
vil lønne sig engang i en tilkommende Evighed.
Men Talen bliver usand, naar det med Eet opdages, og
begge Dele sees i Opdagelsen lige klart: baade, at Selvet,
der forpligtes, er timeligt, og at det forpligtende Gode er
evigt.

Som det levende Væsen er, maa dets Grundkrav,
dets væsenlige Trang ogsaa være. Dyret lever i Timelighed,
det fordrer det Timelige, og kun det Timelige har
Fordringer paa det: det har ingen Forpligtelsens Modsigelse
at drages med. Mennesket, paa hvem der hviler et alvorligt
Krav baade fra det Timelige og det Evige, er stillet
mellem en dobbelt Ild. Ved at holde sig til det Timelige
maa det jo i væsenlig Forstand kunne tilfredsstille sit
inderste Selv, saa sandt dette inderste Selv er timeligt,
kun timeligt. Men saa kommer Evighedstanken med sit
ubetingede Gode og fordrer, at det timelige Selv skal i
dette ubetingede Gode erkjende sit sande Væsen, skal,
skjøndt det veed med sig selv, at det kun er timeligt,
forpligte sig til at leve, som om det var evigt; denne
Modsigelse lader sig ikke løse. Det Væsen, som forpligter,
% --- p. 40 ---
\setcounter{footnote}{0}%
\opage{40}og det Væsen, som forpligtes, maae i deres oprindelige
Grund være af samme Art, i modsat Fald er Forpligtelsen
en Tvangssag uden sædeligt Værd: kun et evigt Selvvæsen
kan forpligtes til det Evige. At fordre af et timeligt Selv,
at det skal lade sig forpligte af et evigt Gode, er lige saa
modsigende, som hvis man vilde tvinge et Dyr, der kun
er organiseret til at bevæge sig i Vandet, til at flyve i
Luften.

Skal det ikke ende med, at de mange Opdagelser tilsidst
gjøre Tilværelsen meningsløs og Livet utaaleligt, da
maa Mennesket endnu en Gang see at opdage sig selv:
det maa opdage Troen og i Troen opdage Aanden.

Men opdage Troen! Her møder den moderne Bevidsthed
med sit Philosophiske. At opdage Troen, siger
man, er endnu mere overflødigt end at opdage Krudtet;
Troen er forlængst opdaget, og dens Hemmelighed gjennemskuet.
Troen er ingen ny Stjerne, men en gammel Hornlygte,
ved hvis magiske Skin Enhver kan see, hvad hans
Indbildning helst vil foregøgle ham. Nei, ikke Tro, men
Viden! Viden er den nye Stjerne, Viden er Verdens Sol,
det kongelige Dagens Lys, der straaler i Tidens Øie. Saaledes
lyder Bebudelsen. Hvad faae vi da at see ved Hjælp
af dette Lys, naar vi, vel at mærke, see ind i Selvet?
Tvivl og Tvivl og atter Tvivl. Selvet i Tvivl og Selvet i
Modsigelse, Evighed i Tanken og Døden i Hjertet: fortærende
Selvmodsigelse! Saa kan Videns Lys i dette Punkt
umulig være det sande Lys. Her er jo desuden ikke eet
Lys, men mange. Lysbølgerne rulle saa vildt, at de mange
Lysskær deelviis forstærke og deelviis igjen formørke hverandre.
De komme ikke fra Aandens Himmel, men fra
Modsigelsernes Skjul, fra uroligt skiftende Tanker. Skiftende
Tanker give kun flagrende Lys: Stjerneskud og Lygtemænd,
Veirlys, som blavre for Tvivlenes Vindpust, slige
vildsomme Lys kunne umulig lede os til at opdage Troen.

Men er Troen da ikke forlængst opdaget? Jo, tilvisse! % til-|visse: the first letters all but lost at the gutter in this copy (z02 Y2); read as the 1884 resetting has it
Men Sagen er, at det med Opdagelser i det per[sonlige] % the word is lost at the gutter in this copy (z02 Y2); supplied from the 1884 resetting
Indre,
i Sindets Dunkelheder, i Selvets forborgne
% --- p. 41 ---
\setcounter{footnote}{0}%
\opage{41}Gruber, forholder sig anderledes end med andre Opdagelser.
Naar En i Videnskaben har gjort en Opdagelse, kunne
Andre vel træde til for at prøve den og, hvis den bestaaer
sin Prøve, tilegne sig den. Men de, som paa denne Maade
tilegne sig den, kunne dog ikke give sig ud for at have
gjort Opdagelsen paa første Haand. Det behøves heller
ikke: det Opdagede staaer fast som et Givet, og det saaledes
Givne kan paa anden, tredie og fjerde Haand komme
til almindelig Kundskab, og dets Erkjendelse udbrede sig i
stedse videre Kredse. Men saaledes forholder det sig ikke
med Opdagelse af Troen og den deraf betingede Opdagelse
af Aanden. Begge hinanden forudsættende Opdagelser ere
Sidebestemmelser af een og samme Grundopdagelse, hvis
Hemmelighed ligger i Selvet, og som desaarsag hvert Menneske
selv maa see at bringe for Lyset. Troens Opdagelse
er en personlig Opgave for hver Enkelt især, altid en Opdagelse
paa første Haand, en primitiv Selvopdagelse. At
denne Opdagelse maa have sine Forudsætninger i den aabenbarede
Religion, er vist nok; men derfor ophører den ikke
at være primitiv. Hvad nytter det vel, at Religionen med dens
hellige Kjendsgjerninger er os given, naar vi ikke kunne
erkjende dens Sandhed. At erkjende Religionens Sandhed
er at opdage Religionen; men at opdage Religionen er
igjen at opdage Troen.

Hvor overfladisk og løs den almindelige Forestilling
om Tro er, giver sig uvilkaarlig tilkjende i tilsvarende Talemaader.
Er det udvortes Ting, historiske Begivenheder,
om hvilke vi endnu ikke have fuld Vished, saa hedder det:
jeg veed det ikke, men jeg troer det. Er det en videnskabelig
Læresætning, hvorom der tvistes, saa hedder det:
jeg vil gjerne troe, at det Opstillede er rigtigt, men jeg
har dog en lille Tvivl; for at jeg kan vide, at Sætningen
er sand, maa jeg have Beviset gjennemført. Saaledes er
det med alle de Sandheder, i hvis Natur det ligger, at de
maae kunne godtgjøres som almeengjældende, være sig ved
umiddelbar Sandsning eller ved Slutninger og tvingende
Beviser. Ved al videnskabelig Erkjendelse gjælder Troen
% --- p. 42 ---
\setcounter{footnote}{0}%
\opage{42}kun for en foreløbig, med Tvivl behæftet Antagelse; almeengyldige
Fornuftsandheder ville \emph{ikke troes}, men \emph{vides}.

Naar nu et tomt, overfladisk Sind eller et af Videnskab
og videnskabelig Forskning optaget og betaget Individ
vil have den samme Maalestok anvendt ogsaa paa det
Religiøse, da er dette et slaaende Beviis for, at den Vidende
endnu ikke har opdaget Troen; thi i religiøs Forstand er
Troen selv den høieste menneskelige Vished. Troen træder
her i Videns Sted og vil desaarsag ikke begrundes paa
Viden. At vide og at troe er i religiøs Forstand det Samme:
jeg veed det, thi jeg \emph{troer} det. Jeg veed, thi jeg troer,
at dette menneskelige Selv, der begynder i Tiden, har et
evigt Ophav og bærer i sig sin Evighedsspire. Jeg veed,
at dette Selv med sin aandig-naturlige Individualitet er
skabt, ikke i Abens Lignelse, men i Guds Billede; jeg
veed det i Tro, saa vist som jeg sandselig veed, at Ilden
brænder, at Sneen kjøler, og Solens Straaler lyse og varme.
Jeg veed det, fordi jeg føler det, og jeg føler det, fordi
jeg tænker det. Jeg veed det, ikke blot fordi jeg føler
det og tænker det, men fordi jeg i alle Følelser, alle Forestillinger
og alle Tanker vil det; jeg vil det, ikke af vilkaarligt
Lune eller blind Indskydelse; jeg vil det, fordi
jeg i rolig, klar, besindig Selverkjendelse maa ville det,
og jeg maa ville det, fordi jeg elsker det. Naar denne
Selvforstaaelse er vunden, saa er Troen opdaget, saa gjælder
det gamle Ord: siig mig, hvad Du elsker, og jeg skal
sige Dig, hvad Du er, thi det, Du elsker, det er Du, og
det, Du væsenlig er, det elsker Du.

Men --- lyder det fra en anden Side --- naar Troen
yder fuldstændig Vished for alle Troende, og Troeslæren
kan jo fornuftigviis kun lægge an paa at have overbevisende
Kraft for de Troende: er det saa ikke heraf indlysende,
at Troens Indhold lige saa fuldt maa kunne afgive
Stof for videnskabelig Begrundelse, som enten Naturens
eller Historiens Indhold; at med andre Ord Troesforskeren
maa staae i samme Forhold til Religionens Kjendsgjerninger,
som Naturforskeren til Naturens? --- Det er en gam%
% --- p. 43 ---
\setcounter{footnote}{0}%
\opage{43}mel og haardnakket, men derfor ikke mindre aandssvækkende
Vildfarelse, vi her have for os, en Vildfarelse, der
er ganske anderledes livsfarlig for Religionens Sag, end
enten Materialismen eller Atheismen eller Nihilismen eller
hvad man ellers vil nævne. Vel maa der gives en Erkjendelseskreds,
i hvilken Religionens Sandheder i deres
indre Sammenhæng og Begrundelse komme til fuld Belysning,
en lys og klar Erkjendelseskreds, i hvilken de evige
Tankelove med et rigt Apparat af gjennemforskede Kjendsgjerninger
vedblive at gjælde. Men i samme Nu man vil
stille Troens Lærebygning paa lige Grundlag med det, som
i egenlig Forstand hører Videnskaben til, styrter hele Bygningen
sammen.

Hvad skal det blive til, naar Troesforskeren vil forholde
sig til Aabenbaringen, som Naturforskeren til Naturen?
Naturen er et altid Nærværende, tilgængelig for
Sandsning og sandselig Vished. Men hvor er der sandselig
Vished for Troen? Hvad skal det blive til, naar
Troesforskeren absolut vil forholde sig til de aabenbarede
Kjendsgjerninger, som Historieforskeren til Historien?

Skulle Troesartiklerne forud begrundes paa historiske
Vidnesbyrd og antages paa udvortes Autoritet, før de kunne
troes, da er Katholicismen --- den pavelige, vel at mærke,
med Dogmet om Pavens Ufeilbarhed --- det eneste Rette;
thi udenfor den katholske Kirke er det dog vel lige saa
umuligt at give fuldt historisk Hjemmel for tre hundrede
Skriftdogmer, som for tre apostoliske Troesartikler. Skal den
personlige Forvisning om Troesartiklernes Sandhed derimod
søges i levende Tro paa det selvnærværende Livets Ord, da er
denne Tro paa Livets Ord en Erkjendelseskilde for sig, et
Princip, der i Selvet begrunder sig selv, uafhængigt af
Viden og Videnskab. Overalt, hvor det formummede Blandingsvæsen
af Tro og Viden endnu fortsætter sit gamle
Spøgeri, er Troen ikke opdaget.

At opdage Troen er at opdage Aanden og i Aanden
at opdage Villien; thi Aanden er Villiens Grund, ja Villien
% --- p. 44 ---
\setcounter{footnote}{0}%
\opage{44}selv. For Menneskevilliens Frihed i Aandens Kraft er Troen
den første Betingelse.

Forskjellen mellem Aandsvillie og Fornuftvillie kan
allerede her blive kjendelig. Fornuftvillien er afmægtig
overfor Naturen, thi det er kun Mennesket, der vil i Fornuftens
Navn, Fornuften selv er intet villende Væsen. De,
der tillægge dette Væsen en fri, selvstændig Villie, maae
enten opstille Fornuften som en Afgud eller paa Grund af
en Sprogforvirring forvexle Fornuften med Aanden. Den
guddommelige Aand har Naturen i sin Magt; Naturen
skabes, bæres og forklares af Aanden. Vil Nogen derimod
paastaae, at det er Fornuften, der har Naturen i sin Magt,
da kunne vi strax vende Sætningen om og sige: det er
Naturen, der har Fornuften i sin Magt. Begge Synsmaader
ere lige sande og lige falske. Det gaaer med Fornuften
som med Lovene: saadanne Love, saadanne Ting;
og omvendt: saadanne Ting, saadanne Love. Tingenes
Væsen udtaler sig i Lovene, men Lovenes Virkelighed beroer
paa Tingene; tage vi Lovene bort, saa forsvinde Tingene,
men dersom vi, i Tanken, tage alle Ting bort, hvad
bliver der saa af Lovene? Det gjensidige Afhængighedsforhold
af Ting og Love, Natur og Fornuft, viser hen til
det tredie Høiere, der har begge i sin Magt, til Aanden.

Aanden har i sin almægtige Villies Eenhed al Naturens
Rigdom og alle Fornuftens Love; den bevæger os
ikke som en dunkel Drift, den hersker ikke som den kolde
Lov; den tvinger ikke, den skaber, den sammenstykker
ikke, den beaander. Det er Aanden, som forklarer det
Individuelle i det Almene, Driften i Fornuft, og Fornuften
i Drift. Aanden vil det Almene, men den vil lige saa
oprindelig det Individuelle. Den opløser ikke Selvet, den
udsletter ikke Skikkelsen; den giver det Evige Liv og
Livet Evighed; den giver det Levende et Selv og Selvet
 Skikkelse, en til Forklarelse indviet Skikkelse: den
opfylder Lovene og den overvinder Døden.

Paa Aandslivets Trin bliver Villien stærk; lad os
see, hvad det er, der giver den Kræfter.

% --- p. 45 ---
\setcounter{footnote}{0}%
\opage{45}Saalænge den menneskelige Villie er henviist til sig
selv alene, saalænge den, saa at sige, for egen Regning
maa føre Kampen mod Naturskrankerne og al den Modstand,
hvormed Nødvendigheden og dens gamle Allierede,
Tilfældene, ere i Stand til at overvælde en indskrænket
Individualitet, seer det fortvivlet ud med Friheden: Drift
og Fornuft staae begge under Skjæbnen. Kan den kæmpende
Villie derimod faae Undsætning af en absolut Styrke,
en Hjælper, som selv er Villie, almægtig Villie, saa klarer
det op i Frihedens Horizont, og Udsigterne blive lyse.
Den almægtige Villie vil ikke fortære den afmægtige, men
optage den i sig og give den Styrke: jo dybere denne
Erkjendelse trænger ind, desto kraftigere maa Villien blive.
Vistnok vedbliver der i Timeligheden at være en Splid
mellem Drifter og Pligter: Driftvillien er selvisk hæftig,
Fornuftvillien tom og kold; men i Aand og Tro --- paa
Aandsvilliens Trin --- er Spliden væsenlig hævet. De
ueensartede Fordringer, de timelige og de evige, begynde
at forvandle sig og blive eensartede. Fornuften er aandig,
den kan beaandes og fyldes med personligt Indhold; Driften
er livsvarm, den kan afkjøles og luttres i Beaandelse,
for igjen at varmes. Thi Aanden er ikke kold; den er
Allivets Væld, i salig Attraa med dybe Længsler, den har
al Inderlighedens Varme.

Gud er Aand, og Gud er Kjærlighed: Aanden er
Kjærlighed. Det ligger i Kjærlighedens Væsen, at den
forener det Modsatte og forliger det Stridige. Fornuften,
det vil sige: vor fornuftigt tænkende Bevidsthed, kan
indsee det Sande i denne Sætning, forsaavidt den fastholdes
i al Almindelighed, men naar Knuden strammes, og det
Afgjørende lægges ind i Kjærlighedens Grund, dens Maal
og dens Villie, bliver Fornuften blind. Skal Fornuften
antage en evig Villie, da bliver det en almindelig, ubevidst
Villen, der falder sammen med den evige Nødvendighed,
med Naturens iboende Drift til at virke og forme. En
guddommelig Villie, som i selvbevidste Øiemed har indblæst
Støvet Liv, har forenet en udødelig Sjæl med et
% --- p. 46 ---
\setcounter{footnote}{0}%
\opage{46}dødeligt Legeme, og spændt Modsætningen af Forkrænkeligt
og Uforkrænkeligt saa stærkt, at dette vort Legeme maa døe,
for at Sjælen i Sandhed kan leve; en Almagtsvillie, som har
gjort alt dette af Kjærlighed og i den faderlige Hensigt,
at Mennesket under Nødvendighedens Tryk og Timelighedens
Byrde skal hæve sig op til Frihed og tilegne sig
et evigt Liv --- en saadan, fri, bevidst og personlig Kjærlighedsvillie
er ikke til for den videnskabelige Fornuft.
Vi maae finde det høist naturligt, at de, som udelukkende
søge al høiere Vished i Viden og Videnskab, vedblivende
ansee en Frihedslære af denne Art for vilkaarlig og ubegrundet.


Fornuften kan indsee, at Frihedsmaalet vilde staae
høit og være et storartet Maal for en kraftig Villies
Stræben, \emph{dersom} der gaves en tilkommende Verden.
Ja \emph{dersom}! Men hvorledes komme vi ud over dette
\emph{Dersom}? Hvorledes bærer man sig ad med at danne
sig et klart Begreb eller endog blot en sandsynlig Forestilling
om en overnaturlig Verdens Beskaffenhed? hvilken
Videnskab paatager sig at føre Beviset for dens absolute
Nødvendighed, uden at indsmugle Forudsætninger, som
ikke tilhøre Viden?

Den Evighedsvillie, som Fornuften kan erkjende, er
i sit uforanderlige Væsen hverken kjærlig eller ukjærlig;
den elsker Intet, thi den hader Intet. Kjærlighed og
Had høre, ligesom Attraa og Afsky, Sympathie og Antipathie,
under Endelighedens Modsætninger, og falde paa Individernes
Affecter.

Ved at mærke os den menneskelige Fornufts Grændse
gjøre vi ikke Fornuften Uret. Tvertimod. At kræve af
et Væsen, hvad der overstiger dets Evne, er ufornuftigt;
det er ufornuftigt at fordre af Fornuften, hvad der er
over dens Kræfter. Den menneskelige Fornuft har gjort
mange vidunderlige Opdagelser i Verden og vil forhaabenlig
endnu gjøre flere; men der er een Opdagelse, den
aldrig har gjort og trods al Gransken, Grublen og Forsken
aldrig vil gjøre, det er Opdagelsen af Selvets evige Liv.
% --- p. 47 ---
\setcounter{footnote}{0}%
\opage{47}Det er just den Opdagelse, der ligger for Troen, den eneste
Opdagelse, som Troen i Aand og Sandhed kan gjøre.
Troen opdager Selvets Evighed: den opdager Aanden og i
Aanden Kjærligheden og i Kjærlighedens Aand et saligt
Liv. Denne Opdagelse er en indre personlig Oplevelse,
den er som Aandedrættet og maa bestandig gjentages;
standser Aandedrættet, saa svinder Livet; standser Oplevelsen,
saa svinder Troen, og med Troen Aanden og med
Aanden den kraftige Villie.

Vi leve i Opdagelsernes Tid; vi leve for at opdage
--- Aandens Villie.

% --- p. 48 ---
\setcounter{footnote}{0}%
\lecture{Femte Forelæsning.}\opage{48}

\emph{Naar Aanden er stærk, saa er Villien} kraftig;
denne Sætning kan indrømmes, uden at den derfor
fører os til dybere Indsigt, thi Spørgsmaalet er: hvorfra
kommer Aanden, og hvorledes vinder den Styrke i
Mennesket?

Det er ikke med Aanden, som med Driften. Hvorfra
kommer Driften? Vi behøve ikke at lede efter den,
den er os nær, i vor Natur. Den \emph{kommer} ikke, den
er, den yttrer sig uvilkaarligt som den er, fordi den er.

Det er ikke med Aanden, som med Fornuften.
Hvorfra kommer Fornuften? Den kommer ikke langveis
fra; vi stirre ikke op i det Blaa og forvente dens Komme.
Den \emph{kommer} ikke, den \emph{er}, i vort Væsen. I samme
Øieblik som vi besinde os paa, hvad det Fornuftige er,
opdage vi, at Fornuften ikke er et Fraværende, men et
os væsenlig Nærværende.

Kan nu det Samme siges om Aanden? Hertil maa
svares baade Ja og Nei. Ja! Aanden er nærværende;
den er jo guddommelig, og det Guddommelige er et Allestedsnærværende.
Og dog Nei! thi Aanden er ikke som
Driften, umiddelbart nærværende i Naturen; den er ikke
for Tanken nødvendigt nærværende, som Fornuften. Naar
det Fornuftige tænkes, saa er det med Nødvendighed; men
fordi En tænker sig noget Aandeligt, deraf følger endnu
% --- p. 49 ---
\setcounter{footnote}{0}%
\opage{49}ikke, at den Tænkende har Aanden selv i levende Nærværelse.
Aanden kan altsaa være paa een Gang baade
nærværende og dog fraværende, kan være os uendelig
nær, skjøndt vi dog maa forvente dens Komme. Det
er som en Gaade. Hemmeligheden er: at Aanden
\textbf{ikke blot er \emph{Væsen}, men ogsaa \emph{Villie}}.

Falder nu Eftertrykket, det for Religion og Tro afgjørende
Eftertryk paa, at Aanden er Villie, saa er dens
Komme en Aabenbarelse, og Aabenbarelsen et Under.
Saaledes kom Evangeliets Aand over Sandhedens første
Vidner; den kom som i en Lyd af Himlen, som i et fremfarende
vældigt Veir, der opfyldte Huset, den kom med
Tunger af Ild og med Sprogenes Gave, den kom med
Evighedens Varsler, med Almagtens Tegn, den kom i et
stort Mirakel. Men dette Mirakel er netop den vidende
Fornuft et Forargelsens Tegn; det er Miraklet, der har uddybet
Skilsmissens Kløft mellem Datid og Nutid. Hvor er
dog ikke Alt forandret siden hiint første Gjennembrud.
Aanden, siger man, er endnu den samme; men hvor er
Kraften, hvor ere dens Undere?

Den civiliserede Nutid troer ikke paa Mirakler, de
Fornuftige kunne ikke engang høre Tale om dem uden at
faae Kvalme. Verden, hedder det, er ikke længer indrettet
paa Mirakler; Mirakler ere dens Constitution imod;
dens Nervesystem fordrager dem ikke, Jordens ikke og
Himlens ikke heller. Vi vende vort Blik nedad, vi grave
i Jorden og finde dernede et stort Bibliothek af Stene,
som i tydelige Skrifttegn berette om forgangne Verdeners
Liv og Tilstande; men i intet Jordstrøg, paa intet Punkt
finde vi Spor af et Paradiis. Vi væbne Øiet og see opad:
Stjerner over Stjerner! men hvor ere Englene, hvor er
den Eviges Throne? Himlen har aabnet sig, ogsaa for
os, men anderledes end i hine Dage. Nutidens Himmel
kan ikke aabne sig, som Datidens aabnede sig for den
Himmelfarende; vore Skyer kunde ikke bære ham, hvis
han engang vilde komme igjen »med Kraft og megen
Herlighed». De Troende see det, og tør dog ikke see
% --- p. 50 ---
\setcounter{footnote}{0}%
\opage{50}det. Troen er en Nonne med Slør, hun slaaer blufærdig
Øiet ned, men stjaalne Blikke gaae gjennem Sløret. Derfor
--- det forsikkre de Skarpsynede --- er Slægten svag,
den er blevet vidende, og vil dog være troende; dens
Væsen er deelt og derfor svækket; den lever halvt i det
Gamle og halvt i det Nye, halvt i overnaturlige Indbildningers
Tryllerier, og dog mere end halvt i Virkelighedens
Interesser med alle en naturlig Verdens nærværende
Kjendsgjerninger: saa sygner Slægten hen i Halvtro og
Tvivl, Villien mattes og Charaktererne svækkes.

Det er ganske vist en alvorlig Krise, der maa
gjennemgaaes, inden den nyere Dannelse faaer sig arbeidet
ud af det Gudernes Tusmørke, hvori dens uforligelige Forestillinger
om Tid og Evighed, lig Jætter og Aser, øiensynlig
stride den sidste Strid. Aabenbaringens Aand med
sine Undere og Naturen med sine uforanderlige Love have
i vor Nutidsbevidsthed stillet sig saa skarpt mod hinanden,
at Halvtroen gaaer med en ond Samvittighed og venter
paa en sidste Forklaring. Troestanker og Videnstanker i
forvirret Blanding, Phantasier af Religion og Indtryk af
Videnskab bevæge sig i underlige Syner, vildledende, som
Ørkenens Luftspeilinger for den Reisende. Veien gjennem
dette Uføre, den eneste Vei, som sikkert fører til Maalet,
er betegnet ved et gjennemgaaende Spørgsmaal, Spørgsmaalet
om Villien. Dette Spørgsmaal har nemlig det
Særegne ved sig, at det fører Striden om alt det Udvortes
tilbage til et eneste, afgjørende Indvortes.

Hvad er Altilværelsens inderste Bevægende: Naturens
dunkle Drift eller en evig fri, almægtig Villie? Det er Cardinalpunktet,
hvorom det Hele dreier sig. Antages det
Første, saa er Miraklet en Umulighed, og historiske Vidnesbyrd,
de være nok saa mange og nok saa overeensstemmende,
kunne ikke blæse Virkelighed ind i det Umulige.
Antages det Sidste, saa er Miraklet muligt. Og dog kunne
historiske Vidnesbyrd ikke give os fuld Vished om nogetsomhelst
enkelt virkeligt Mirakel. Historiske Vidnesbyrd
give os kun en Sandsynlighed; men det virkelige Mirakel,
% --- p. 51 ---
\setcounter{footnote}{0}%
\opage{51}det ligger i dets Væsen, er altid usandsynligt. Viden holder
sig til det historisk Sandsynlige, og forkaster det Usandsynlige;
Troen gaaer ud over det Sandsynlige, og griber
det Mulige. Det Umulige kan ikke troes, men det Usandsynlige,
som dog er muligt, kan troes: her er Forskjellen
mellem Tro og Viden.

Paa Profanhistoriens Omraade har den forskende
Videnskab en sikker Veiledning i bekjendte Love og en
derved normeret Sandsynlighedens Maalestok; naar en Tildragelse
er saa sandsynlig, at Ingen kan negte den uden
med det Samme at benegte al historisk Sandsynlighed, saa
er Tildragelsen sikkret som Kjendsgjerning. Anderledes
med Aabenbaringens Kjendsgjerninger i den hellige Historie.
Her kan det Sandsynlige ikke normeres efter bekjendte
Love, thi Underet selv er det Usandsynligste af
Alt. Griber nu Troen ud over det blot Sandsynlige for, i
den almægtige Villies Navn, under Virkelighedens Form
at gribe den høieste Mulighed, saa synes den jo udsat
for at oversee Tingenes naturnødvendige Sammenhæng,
glemme Lovene og komme ind paa det Eventyrlige. Imod
denne Vildfarelse staaer Livets Ord, det guddommelige Selv,
Magtvilliens Aabenbarer, fast, urokkelig fast som Klippen
i Havet. Hans Villie er Magt, alle Muligheders Hemmelighed,
han er Underet selv, omgiven af Under og Tegn.
Og dog er hans egen Villie, menneskelig seet, som en
Afmagt; han har, bogstavelig talt, ingen Egenvillie: hans
Villie er at gjøre Almagtens Villie, Faderens Villie, som
er i Himlene. Saa staaer han paa Frihedens Grund, og
dog under strenge Love. Han viser sig, som han er,
--- Sandheden! Ja, men i Lydighed. Ikke som en vilkaarlig
Hersker, men som en lydig Tjener, lever han i Sandheden,
strider for Sandheden, lider for Sandheden, døer
for Sandheden: han, den tjenende Sandhed, bryder ikke
Lovene, han opfylder Lovene. Han stifter ikke en ny
Religion ved Hjælp af en ny Lære med et godt Exempel,
nei, han stifter, hvad han er, og han er, hvad han
stifter, Religionen selv, \emph{Aandsvillien} i Kjød og Blod.

% --- p. 52 ---
\setcounter{footnote}{0}%
\opage{52}Det gaaer jo feiende nok med Erklæringer som denne:
»vi leve i Lovenes, ikke i Miraklernes Tid»; men er saa
Menneskeslægtens høieste Anliggende hermed afgjort?
Lad os oversætte Erklæringen i et andet Sprog, for ret
at see, hvad den indeholder. Vi leve i Lovenes Tid, det
vil sige: Almagtsvillien er gaaet ud, kun de endelige
Villier ere tilbage. Vor Fremtids Fane bærer den nye
Indskrift: \emph{Døgnvillien}. Hermed stemmer det nu ganske
vel, at alle Døgnvillier skulle frem, og Fleertalsvillien
sidde paa Magtens Throne; men hvorledes denne Døgnvillie
vil bære sig ad med at naae det høie Maal: at opfylde
Lovene og løfte Masserne op paa Frihedens Trin, er endnu
ikke forklaret.

Spørgsmaalet om Underet falder i sidste Instans
sammen med Spørgsmaalet om \emph{Villiens Oprindelighed}.
Dette Spørgsmaal lader sig ikke affærdige ved at paaberaabe
sig Naturlovenes Uforanderlighed. Man kunde da
lige saa godt henvise til Tankelovenes Uforanderlighed.
Kunne de uforanderlige Tankelove ikke engang siges at
gjøre den menneskelige Villie Afbræk, naar Talen er om
nye Ideer, nye Begyndelser og tilsvarende nye Handlinger,
hvor meget mindre skulle de uforanderlige Naturlove da i
noget Punkt kunne gjøre den guddommelige Villies ubetingede
Frihed Afbræk. Den ubetingede Frihed bestaaer
ikke deri, at Villien er uden Betingelser, men deri, at den
har Betingelserne i sin Magt, og er hævet over endelige
Skranker. Fornuften er Formen for Viisdommens Indhold,
men Indholdet selv bestemmes af Villien. Aanden har
ikke Fornuften udenfor sig, saa lidt som den har Viisdommen
udenfor sig. Den almægtige Villie kan gjøre,
hvad den vil, virke og røre sig, som den vil, den kommer
aldrig udenfor det uendelige Heles uendelige Sammenhæng:
den kan jo aldrig komme udenfor sig selv. Men netop
i Sammenhængens Uendelighed har Friheden sine uudtømmelige
Muligheder.

For den almægtige Villie ere alle Ting mulige; det
betyder ikke, at Villiens Almagt kan gjøre Lovene til Intet,
% --- p. 53 ---
\setcounter{footnote}{0}%
\opage{53}ophøie Modsigelser til Love, gjøre Natur til Unatur, Cirkelen
fiirkantet o. s. v., men det betyder, at den i og ved
og med Lovene kan bringe det i Menneskenes Øine Allerusandsynligste
frem, kan uden Brud paa Tingenes skjulte
Sammenhæng pludselig bryde alle Sandsynligheders Maalestokke.
Den almægtige Villie, der har Tingene i sin Magt,
har ogsaa Lovene i sin Magt; men at have Lovene i sin
Magt er: at have deres Opfyldelse i sin Magt. Den, der
har Tanken i sin Magt, gjør aldrig Brud paa Tankelovene;
den, der har Naturen i sin Magt, gjør aldrig Brud paa
Naturlovene.

Der er en Grundsætning, hvorom Aabenbaringstroen
med sin Forvisning om Aand og Villie maa være enig,
aabenlyst, oprigtigt, fuldstændigt enig med Nutidsbevidstheden,
det er Sætningen om Lovenes Uforanderlighed.
Naar denne Sætnings Gyldighed, ret forstaaet, tilfulde er
anerkjendt, saa er Forbindelsen mellem Fortid og Nutid
aandelig gjenoprettet, og Muligheden af en gjensidig Forstaaelse
mellem Fornuft og Religion tilveiebragt. Paa
Grundlaget af denne Forstaaelse beroer Adskillelsen af det,
som videnskabelig kan vides, og det, som ikke kan vides,
men maa falde ind under Tro. Hvad er den oprindelige
Grund til Naturlovenes Uforanderlighed? »Tingenes Væsen,
Verdens fornuftnødvendige Sammenhæng»; det er Naturlærens
Forklaring. »Guds almægtige Villie» ; det er Religionens
Forklaring. Her viser sig paa Erkjendelsernes
Landkaart en tydelig Grændse mellem Tro og Viden, en
Demarcationslinie, nok saa tilforladelig som den, Paven
drog, da Columbus havde opdaget den nye Verden.

Erfaringslæren begynder med Kjendsgjerningerne,
holder sig til Lovene og standser ved Nødvendigheden.
Naturphilosophien gaaer et Skridt videre. Den kan i Grundforholdet
af Natur og Aand paavise den betingede Nødvendigheds
Fremgang af Friheden, Lovenes Oprindelse fra den lovgivende
Villie. Men den maa med alt dette dog blive
staaende ved Nødvendighedssiden og nøies med det, som
uafviselig hører med til Forstaaelse af \emph{givne} Kjends%
% --- p. 54 ---
\setcounter{footnote}{0}%
\opage{54}gjerninger og videnskabelig Begrundelse af \emph{bekjendte}
Love. Som  i en hellig Skumring af halvgjennemsigtig
Nødvendighed kan Naturphilosophien overalt see Lysskæret
af den guddommelige Frihed; men Villien selv i
Mulighedernes indre Uendelighed kan den ikke see: Frihedens
Sol gaaer aldrig op i disse Videns Egne, den
holder sig i Horizontens Rand.

Naturphilosophiens sidste Sætning er Religionens
første: Naturen bæres af Aanden, Lovene grunde i Villien.
Nu gaaer Frihedssolen op, det dages over Mulighederne.
Det er en Erkjendelseskreds for sig, Troens, Aabenbaringens
evige Lyskreds. Hvad see vi her i dette nye Lys?
Villien i det almægtige Selv, Friheden i alle sine Muligheder:
\emph{Underet}! Underet er, at Almagten er Villie, og
Villien Almagt; Underet er, at Alverdens Nødvendighed
har sin oprindelige Grund i guddommelig Frihed, Underet
selv er Viisdommens Hemmelighed.

Det Under, hvorom her tales, er ikke Phænomen,
men Væsen, det er et Grundunder. Det har sit ubedragelige
Vishedsmærke, ikke i udvortes Tegn, men i indvortes,
i vort personlige Selv, i Selvforstaaelse af Tro og Villie.
Det skal ikke først stadfæstes ved vidtløftige Vidneforhør
over ydre Mirakler; det er omvendt de ydre Mirakler, som,
naar Troeskritiken anlægger sin Maalestok, alle maae føres
tilbage til og stadfæstes af det altid værende, evigt blivende,
aldrig skiftende Grundunder.

Naar vi nu, under disse Forudsætninger, gjentage, at
Aanden er Villie, og dens Aabenbarlse % PRINTED AS IS (Aabenbarelse)
et Under: saa er
det dog vel indlysende, at heri ikke skjuler sig mindste
Bagtanke eller Bitanke om Lovbrud, løs Vilkaarlighed eller
guddommeligt Lune. Aanden er Væsen, og Aanden er
Villie. At Aanden aabenbarer sig, vil da sige: den viser
sig i Troens Speil som den er efter sin Villies evige
Væsen. Hvad er da Aanden efter sit aabenbarede Væsen?
Kjærlighed, hellig, uforanderlig Kjærlighed. Heri ligger,
at den altid er redebon --- thi Kjærlighed er altid rede%
% --- p. 55 ---
\setcounter{footnote}{0}%
\opage{55}bon ---
at den altid vil komme, naar og hvor den, paa
Frihedens Vilkaar, kan komme.

Vi kjende Sandhedens Aand paa Livets Ord. Spørge
vi, hvorfor den da tøvede saalænge, inden den kom ---
med Livets Ord til fuld Aabenbarelse, da er det Troens
eenfoldige Svar: den kunde ikke komme før, den var forhindret.
Den var i Ordet, og Ordet var Verdens Lys, og
Lyset skinnede --- men i mørke Myther, og Mørke og
Myther begrebe det ikke. Spørge vi, hvorfor den ikke
kommer til os i Magtens Tegn med Underets Kraft, som
den kom til Fædrene, da er det Troens eenfoldige Svar:
den kommer ikke med Underets Kraft, fordi den ikke kan,
den er forhindret. Det er den menneskelige Villie, som
paa Frihedens Vei har kastet de mange Stene og opdynget
de mange Hindringer.

»Men», siger man, »saa er her jo dog en Modsigelse,
en utaalelig, fra Under og Villie uadskillelig Modsigelse.
Eet af To: enten er Aanden almægtig, eller den er afmægtig.
Er Aanden almægtig, hvad kan da standse dens
Komme? Den kommer og gaaer, som den vil, naar den
vil, og hvorhen den vil; den kommer med Seier, lysende,
varmende: Hjerterne smelte for Ord af Ild, som Klippernes
Iis for Solens Glød. Saaledes kommer en Almægtig. Kan
Aanden ikke komme saaledes; er der Noget, som kan
stille sig i Veien, et virkeligt Noget, som kan lamme
dens Kraft, tvinge den til at staae stille eller dog bøie
af og liste sig frem ad Omveie: hvor er saa Villiens Almagt,
og hvad er Underet?»  Unegteligt! her er en Modsigelse.
Gudsvillien er almægtig, Menneskets Villie er afmægtig.
»Og saa skulde dette Afmægtige være i Stand
til at sætte sig op mod det Almægtige, saa skulde det
Selviske i Menneskets Egenvillie virkelig have Kraft til at
gjøre Aandsvillien Modstand? Det maa være en Indbildning,
et Bedrag». Vi svare: Ja! det er ganske vist en
Indbildning, at Døgnvillier, om de end i Millionviis sammenrottede
sig til Oprør, skulde formaae at sætte den evige
Villie Skranker; det er et Bedrag, et gruopvækkende,
% --- p. 56 ---
\setcounter{footnote}{0}%
\opage{56}aandsfortærende Selvbedrag. Men dette Selvbedrag er et
virkeligt Noget, det er i Menneskets Egenvillie saa virkeligt,
at det kan »bedrøve Aanden».

Kom det an paa Magten alene, saa var Aanden vis
paa sin Seier: det ringeste Pust, og Alverdens Hindringer
vege som Fnug. Men hvad blev det saa til med Menneskets
Frihed? Skulde Aanden bruge Magt og tage den
gjenstridige Egenvillie med Vold, saa maatte den jo grundlægge
Friheden ved at gjøre Friheden til Intet: det er en
Modsigelse saa klar, at vi heri see det Umulige. Var
Talen om at gjøre de Døve hørende, de Blinde seende, at
aabne Gravene og opvække de Døde, saa kunde vi ikke
sige: det er Almagtsvillien umuligt; thi vor Viden slaaer
ikke til. Vi kunne sige, at slige Gjerninger ikke lade sig
forklare efter de os bekjendte Love; men vi kunne ikke
sige, at de staae i Strid med Lovene, thi de endnu uopdagede
Loves Tal er uendelig større end Tallet paa de
opdagede. Alle de Modsigelser, som en letfærdig Viden i
sin Overilelse i denne Henseende har formuleret, ere byggede
paa Sand og falde bort af sig selv, naar de prøves. Men
den Modsigelse, at den almægtige Villie skulde ikke blot
kunne tvinge den afmægtige --- deri er ingen Modsigelse
--- men tvinge den til ikke at lade sig tvinge, tvinge den til
at bestemme sig med Frihed, det er en Modsigelse saa klar
og indlysende, at den, der tillægger Almagten en saadan
Mulighed, ikke veed, hvad han selv siger. Frihedens Knude
vil ikke hugges over, den vil løses.

Knudens Løsning er Selvbedragets Opløsning. Selvbedraget
er i Mennesket dobbelt, ja tredobbelt. Det er i
Tanken og kaldes Forblindelse; det er i Villien og kaldes
Trods; det er i Følelsens Eenhed af Tanke og Villie, og
kaldes Forhærdelse. Det er vanskeligt, om ikke umuligt,
at komme Bedraget tillivs under nogen enkelt af disse
Former for sig; de tre Medinteressenter, Tanke, Villie og
Følelse, holde sammen i det fælles Bedrag, de ere som
Sammensvorne: den Ene vil ikke forraade den Anden.
Forstanden kan ikke opklares, og Forblindelsen ikke hæves,
% --- p. 57 ---
\setcounter{footnote}{0}%
\opage{57}saalænge Villien er vrang og Følelsernes Strøm gaaer
Sandheden imod. Følelsen faaer falske Indtryk, saalænge
Forstanden er omtaaget, og den af falske Forestillinger
bedaarede Villie holder Indgangen spærret for enhver Lysstraale.
Og skulde vel en Villie med falske Tanker og
falske Forestillinger, falske Billeder og selviske Indskydelser
paa egen Haand være i Stand til at begynde med Sandheden
ud fra sig selv og hugge Bedraget igjennem, som en
Chlotaldo sønderhugger Morganas Perle? Det Hele viser
tilbage til Selvets uopladte Naturgrund, den ubevidste Eenhed,
som bærer Tanke, Villie og Følelse i uforklaret Selvbestemmen:
skal Bedraget kunne hæves, saa maa der i
denne Naturgrund være Noget, som holder igjen og
gjemmer en Mulighed for Aanden.

Kunde vi forestille os Selvbedraget fortsætte sit Kredsløb
gjennem de tre Hovedformer af Tanke, Villie og Følelse
saa længe, saa vedholdende og i en saadan Dybde, at det
Skjulte i Selvets Grund blev bragt op i Bevidsthedens Lys
og den sidste Rest af naturlig Umiddelbarhed opbrugt:
saa vilde Selvbedraget fortætte % PRINTED AS IS (fortsætte)
sig i en Villie af ubetinget
Lysfornegtelse, en i Aand formørket Villie, en Aandsvillie,
der ikke var menneskelig, men satanisk. At det under
denne Forudsætning hverken vilde være Tanken eller Følelsen,
men Villien selv, der til allersidst kom til at gjøre
Udslaget, er indlysende. Følelsen er uvilkaarlig; det
følende Væsen har ikke Indtrykket i sin Magt, og heri
ligger altid en Undskyldning for Følelsesbedraget. Ogsaa
Tankebedraget har sin Undskyldning, saa længe det Ubevidste
endnu er saa fremherskende, at der kan bevares et
Skin af, at Tanken bestandig er i Begreb med at søge
det Sande. Først naar den fornegtede Sandhed staaer
klaret for Tanken og Smertens Stik gaaer igjennem Hjertet,
har Egenvillien Betingelser for at forvandle Selvbedraget
til absolut Selvforblindelse og drive Forhærdelsen til det
Yderste.

Dersom en saadan Selvmodsigelse blot var logisk,
vilde den uden videre være at henføre under det Umulige,
% --- p. 58 ---
\setcounter{footnote}{0}%
\opage{58}og Bedraget glide Tanken forbi, som en Skygge paa Væggen.
Men en Selvmodsigelse af denne Art er mere end logisk,
den er personlig. Allerede paa Driftvilliens Trin gjør
Individet Erfaringer af, at det Sande ingenlunde er det
Samme som det Behagelige, og Forstanden faaer ved slige
Anledninger ofte travlt nok med at tilhylle det Usande
og besmykke det Halvsande. Paa Fornufterkjendelsens
Trin kan Egenvillien i Skjæbnens Navn holde sig selv saa
fast, at den hellere vil »knuses» end »bøies». Kommer
nu en saadan potenseret Egenvillie i Forhold til Aanden
selv, Sandhedens og Kjærlighedens Aand, da er det ingenlunde
givet, at den bøier sig for Aanden, om den end
maa erkjende Aanden for Magt. Det Modsatte kan skee.
Den potenserede Egenvillie føler, at den i al sin Afmagt
dog har en Magt, Magt til at trodse Aanden. Greben af
Selvmodsigelsens Krampe, griber den da med Fortvivlelsens
Kraft til det Valg: hellere at lide alle Selvfortærelsens
Kvaler end bekjende sin Skyld, bære Ydmygelsen og bøie
sig for den Overmægtige. Naar dette Valg er skeet, saa
er Forblindelsen fuldendt; saa er Sandhed Løgn, og Løgn er
Sandhed, Frihed Trældom, og Trældom Frihed: Kjærlighedens
Lys udslukkes i Hadets Mørke.

Den samme Sandhedens Aand, der vækker Kjærlighed,
kan ogsaa vække Had, thi Had er en omvendt, i selvisk
Trods forvrænget Kjærlighed. Det er ikke Kjærlighedens
Skyld, at den vækker Had; den kan ikke gjøre for det.
Kjærligheden kan taale Alt for Andres Skyld, men Selvmodsigelse
kan den ikke taale. Den kan ikke fornegte
sit Væsen og sige til sig selv: jeg maa forandre mig, jeg
maa see at blive egenkjærlig, det vil sige: ukjærlig, thi
ellers vækker jeg Had. Kjærlighedens Væsen er: ikke at
skjule, men at aabenbare sig. Naar da den sig aabenbarende
Kjærligheds Aand vækker Had, saa bliver Hadet
ikke ringe, fordi Kjærligheden er stor; nei, Hadet nærer
sig af Kjærligheden, voxer ved Kjærligheden, stræber at
maale sig med Kjærligheden og blive uendeligt, fordi
Kjærligheden er uendelig. Hvad heri er skjult, skulle vi i
% --- p. 59 ---
\setcounter{footnote}{0}%
\opage{59}denne Sammenhæng ikke forsøge at drage frem, men kun
holde det fast, at der paa Villiens Vei i Frihedens Navn
virkelig er et Noget, som kan bedrøve Aanden.

Vi have betragtet Selvbedraget fra den stærke Side;
lad os til Slutning vende Synspunktet om og betragte det
fra Svaghedens Side.

Der er jo Mennesker, og deres Tal er vist ikke ringe,
bekymrede Sjæle, som sukke i sig selv og sige: «Vi staae
ikke Aanden imod; vi føle, at vi trænge til Aanden; vi
bede om Aanden, men --- Aanden kommer ikke.« Og der
er til Vederlag igjen dem, som i Overmaal nyde Aanden,
Troens Udvalgte, hvis fedtpandede Trivelighed ligefrem glindser
af Aanden: de snakke om Aanden og snakke med Aanden.
De besidde det Trylleord, hvori Aanden boer; saa kjende
de Aanden godt: de ere, reent ud sagt, Kammerater med
Aanden. Der er i Forholdet mellem Gud og hvert
enkelt Menneske en Hemmelighed, som kun den Alvidende gjennemskuer; men see vi lidt skarpt paa det Menneskelige,
viser det sig da ikke ogsaa her, at Selvbedraget er
dobbelt, ja tredobbelt.

Der er et Forestillingsbedrag, som med Forglemmelse
af Aandens Allestedsnærværelse, tager Aabenbaringsordet
forfængeligt, sætter Aanden fast i et Billede og gjør Billedet
udvortes. Under Billedet af en usynlig, men dog
endelig Person, der bevæger sig i Rummet, farer »den
Helligaand«, op og ned og ned og op, et bevinget Mellembud
mellem Jord og Himmel. Vi forarges ikke over det Billedlige;
men naar dette Billedlige krystalliserer i Phantasien
til en fast Virkelighed, naar Aanden til visse Tider
aflægger Besøg i et Menneskehjerte og gaaer igjen med
Løfte om at vende tilbage; naar, sige vi, dette Billedvæsen
faaer en saadan Magt med Bevidstheden, at den Bedende
sidder og venter, at Aanden som Gjæst skal banke paa
Døren: tør vi da undre os, om en saadan Ventende kan
komme til at vente forgjæves? At gjøre Personlighedens
Aand til et mythisk Væsen, trolovet med Iris, der farer paa
Regnbuen, er at opfatte den som en Afgud; men at opfatte
Aanden som en Afgud, er at bedrøve Aanden.

% --- p. 60 ---
\setcounter{footnote}{0}%
\opage{60}Det antydede Forestillingsbedrag forbinder sig med
et tilsvarende Følelsebedrag. % PRINTED AS IS (Følelsesbedrag)
 Indbildningen er, at Aandens
Nærværelse skulde være Gjenstand for en umiddelbar indvortes
Sandsning, at den Sjæl, hvem Aanden besøger,
skulde være i Stand til at sige paa Klokkeslet, naar Besøgelsen
skete, skulde kunne gjøre en tydelig Adskillelse
mellem to Væsener, sit eget Væsen og Aanden,
skulde ligefrem kunne samtale med Aanden: »det sagde
jeg til Aanden, og saa sagde Aanden til mig« o. s. v.
Naar denne Indbildning faaer Overhaand, og dens skjulte
Indflydelse er uden Tvivl langt større end Mange ret
ville tilstaae: tør vi da finde det forunderligt, at der er
oprigtige Sjæle, som vedgaae for sig selv, at de aldrig
have staaet Ansigt til Ansigt med Aanden selv, aldrig
nogensinde have enten seet eller umiddelbart fornummet
Guds egen Nærværelse som Helligaand?

Saa skal den hellige Skrift bære Skylden. I Modsætning
til de Billedtroende, som, for at tage Skriftordet ret
bogstaveligt, gjøre det Overnaturlige unaturligt, og sætte det
Indvortes til paa et Udvortes, optræde saa »de Fornuftige«
med en Vidensfortolkning, som bortskaffer det Overnaturlige,
udvander Aabenbaringsordet og faaer dets Charakteerpræg
udvisket. Saa bliver Aandens Bog ved Hjælp af et
dobbelt Selvbedrag til en tvetydig Bog: Børnene i Ammestuen
holde sig til Billederne og lege Juul med det Overnaturlige;
de Voxne opdage Naturligheden og lefle med
Mythologien.

Forestillings- og Følelsesbedraget forene sig i et Tredie,
den farligste Form af dem alle, og det er Villiesbedraget.
Det Skuffende, hvorfra man gaaer ud, er Forestillingen om,
at Menneskets naturlige Villie kan blive tilovers som en
Størrelse for sig ved Siden af Aandsvillien. Hos Nogle bliver
dette opfattet saaledes, at den naturlige Villie, være sig i
Form af Driftvillie eller Fornuftvillie, skal tvinges ind
under Aandens Lov og pidskes op til, om ikke just at fortjene
Aanden, saa dog til at gjøre sig værdig til Aanden,
arbeide sammen med Aanden, og, om muligt, ved gode
% --- p. 61 ---
\setcounter{footnote}{0}%
\opage{61}Gjerninger indlægge sig Fortjenester hos Aanden. Herhen
hører en egen Form af den saakaldte strenge, bekymrede
Christendom, ikke den i Sandhed, men den i Suurhed bekymrede,
en Christendom, hvis Kjærlighed er en Art
gudelig Ærgrelse, hvis Trøst er en Forvisning om de Glades
evige Fordømmelse, en Christendom, som hader Verden med
en guul Tro, der har sit Sæde i Galden. Hos Andre bliver
Menneskets naturlige Villie ganske rigtigt erklæret for
uduelig til at gjøre Aandens Gjerninger. Alt beroer uden
videre paa Aanden og Troen, den nye Villie kommer af
sig selv. Vil Aanden, saa er der Villie; vil Aanden ikke,
nu vel, saa er der endnu ingen Villie --- Villien groer
langsomt. Men det gjør Ingenting: af Tro er der fuldt op,
og Aand og Tro er det Hele. Paa disse Betingelser
kommer det til et skjønt Forlig mellem Aandsmennesket
og det naturlige Menneske: de kunne nu saa godt sammen,
thi de elske jo begge Friheden. Troen har sin Frihed
i det Overnaturlige, og det kan jo ikke genere det
Naturlige; Driftvillien, der har sin Frihed i det Naturlige,
kan gaae omkring og græsse, hvor helst den vil, paa Timelighedens
Overdrev, naar den for Resten er skikkelig og
ikke forgriber sig paa Troen. Forliget er saa hjerteligt, at
det livagtigt seer ud, som om de to Villier, Troens og
Driftens, vare blevne til een Villie. Herhen hører en vis
Formation af den saakaldte glade Christendom, ikke den i
 Sandhed glade --- thi det er for Troen evig vist, at Aandens
Liv er et Liv i Glæde --- men den i Naturlighed
glade, den kisteglade, der finder det saa deiligt paa Jorden,
fordi den er saa sikker paa Himlen.

Naar disse Selvbedrag hæves, og de kunne hæves
--- hvert Menneske har i det aabenbarede Ord, det Ord, hvis
levende, sakramentale Nærværelse er at finde i Menigheden,
Betingelser for at hæve dem, naar han alvorlig vil --- saa
kommer Aanden.

Der gjøres i vor nærværende Tid adskjellige Forsøg
paa at forlige Nutidsbevidstheden med den bestaaende Religion.
Man gaaer i denne Henseende ud fra den ganske
% --- p. 62 ---
\setcounter{footnote}{0}%
\opage{62}vist almeengyldige Forudsætning,  at Menneskeheden har
Trang til Religion, en saa dybtgaaende og væsenlig Trang,
at Stater og Folk, Civilisation og Sædelighed, al borgerlig
Orden og alle en høiere Oplysnings Goder vilde løse sig op
i et Chaos af raae Lidenskaber og Menneskelighedens Træk
forvrænges til Rovdyrets Glubskhed og til Abens Lignelse,
hvis Religionen forsvandt fra det menneskelige Samfund.
Man holder endvidere fast paa den i sig fuldgyldige
Sandhed, at det Religiøse umuligt kan være os et Levende,
Nærværende, dersom det ikke paa aandskraftig Maade formaaer
at gjennemtrænge Nutidsbevidstheden saa fuldstændig,
at vi med fuld Frihed leve og virke og røre os deri
som i vort rette, aandige Livselement. Man har ganske
rigtig indseet, at det maa være dette aandige Livselement,
der indeholder de første afgjørende Betingelser for Dannelsen
af Charakterer med et nyt Syn og en ny Villie,
en Villie, som med Tro og Mod tør gribe Opgaverne fra
Grunden af og freidig arbeide paa deres Løsning. Om alt
dette maae vistnok de Fremmeligere i Tiden alle være
enige; men, gaae vi nærmere ind paa Sagen, møde
vi saa ikke dette Forunderlige, at Forsøgene begynde med
at oversee det første Nødvendige?

Skal det Religiøse sættes i Kraft, saa er det dog vel
aldeles nødvendigt, at Enhver maa begynde med at spørge
sig selv, hvori det Religiøse da væsenlig bestaaer, hvad
Religion oprindelig er, og hvor Sandhedsmærket findes,
det ubedragelige, der kan forvisse ham om, at han ikke
farer vild og begiver sig ind paa det Uvisse. Er Sandhedsmærket
tilforladeligt, da maa det findes i hans eget
Selv, i Selvforstaaelse af Tro og Aand. Denne Selvforstaaelse
er en Villiesforstaaelse, en Forstaaelse af Underet;
thi Underet er Religionen selv. At en menneskelig Villie
staaer i et Nødvendighedsforhold til villieløse Ting, deri
er der intet Under; at menneskelige Villier forholde sig til
hverandre som Vilkaarlighedspunkter, der alle beherskes af
overmægtige Love, deri er der intet Under. Men at Guds
almægtige Villie er traadt i Forhold til Menneskets afmæg%
% --- p. 63 ---
\setcounter{footnote}{0}%
\opage{63}tige Villie, har aabenbaret sig selv som en Menneskevillie,
og ved denne sin Selvaabenbarelse stiftet et Kjærlighedsforhold,
en salig Forligelse, at de to Villier, Guds og Menneskets,
til evig Tid kunne være som een Villie, det er i
Aand og Sandhed et Under. De, der mene, at de bedst
kunne sætte det Religiøse i Kraft ved at udelukke Underet,
ere som de, der mene, at de bedst kunne tegne Cirkelen
ved at fremstille den som en Fiirkant. At forkaste Underet
er at forkaste Religionen; at antage Religionen som Religion
er at erkjende Underet.

Staaer Sandhedsmærket fast i det Indre, saa udgaaer
herfra det Lys, som kan opklare det tilsvarende Ydre.
Det Indre i Troen er saa langt fra at være ligegyldigt mod
sit Ydre, at hver Troende maa erkjende, at hans Selvforstaaelse
i Aand ikke skyldes ham selv, men er vakt ved
et udenfra kommende Ord, og hviler paa det sakramentale
Ordets Under, som med Aandsvilliens Almagt meddeler sit
Indre i, med og under et bestemt Ydre. Saa er her
altsaa formedelst Underet et fortsat, levende Vexelforhold
mellem det Indre og det Ydre. Saalænge alt det Ydre
opfattes og bedømmes som Middel og Gjennemgangsled for
det Indre, saalænge vedbliver Selvforstaaelsen at berige sig,
og Sandhedsmærket slaaer til. Først naar Bevidstheden
gjør Forsøg paa at gaae uden om Troen, bagved Troen,
for at gribe den religiøse Kjendsgjerning i reen, historisk
Udvorteshed, svigter Sandhedsmærket, og Famlen og Tvivlen
begynder. Det kan ikke noksom indskærpes, at Troen og
Troens Indhold ere uadskillelige: ethvert Mirakel er Karikatur,
naar det ikke bliver seet med Troens Øie. Den
hellige Historie udfolder for vor Betragtning en Række
Livsbilleder og fører os ind i en Kreds af Kjendsgjerninger
og Phænomener, i hvilke Fortidens Troende, de være nu
Øienvidner eller ikke Øienvidner, have seet og erkjendt de
Almagtstegn, ved hvilke den frelsende, forløsende Kjærlighedsvillie
har aabenbaret sig for dem; Spørgsmaalet er,
om vi i deres Tro kunne gjenkjende Troens Væsen og dele
Tro med dem. Forstaaelsen og Tilegnelsen af den hellige
% --- p. 64 ---
\setcounter{footnote}{0}%
\opage{64}Historie er en sig fortsættende Gjenkjendelse af Troens
Liv, dens Væsen og Indhold. Hvor denne Gjenkjendelse
svigter, der svigter ogsaa vort Sandhedsmærke, og det
Overleverede berører os ikke religiøst.

Det er en fuldt berettiget Fordring, den nyere Tid
har opstillet, at Forskningen skal være fri. Paa de her
opstillede Vilkaar kan Aabenbaringstroen ikke have det
Ringeste at indvende imod den videnskabelige Forsknings
ubegrændsede Frihed; thi nu veed den med sig selv, at
ingen Videnskritik, dens Resultater være nok saa negative,
længer er i Stand til at undergrave det religiøse Grundlag.
Blev det end videnskabeligt beviist, at vi ikke havde
et eneste Evangelium umiddelbart fra noget Øienvidne, ja
selv om det kunde bevises, at vort Matthæusevangelium
var fra det andet, og vort Johannesevangelium fra det
tredie Aarhundrede: saa vilde alle disse Beviser dog ikke
have den ringeste Indflydelse paa Bestemmelsen af Religionens
Indhold. Her spørges ikke om de historiske
Mellemled, hvorigjennem Sandhedsmeddelelsen er gaaet;
her spørges om Sandheden, den aabenbarede Sandhed, og
Sandhedsmærket, det ubedragelige, er i vort Indre.

Vi leve af Underet og dog under Lovene; vi leve
under Nødvendighedens Tryk og dog i Friheden, midt i
den materielle Verdens haandgribelige Kjendsgjerninger og
dog i en nærværende Evighed, saa sandt vi have den personlige
Vished, at Villien er Aand, og Aanden Villie.

At maale sin Tro med Følelsens Varme er usikkert,
thi Varmemaalet er ubestemmeligt, og det Ubestemmelige
kan skuffe. At maale sin Tro med Tankens Klarhed er
usikkert, thi Klarheden er kold og Idealet i Tankens Speil
staaer udenfor Virkeligheden. Men at maale sin Tro med
Aand og Villie, med den Aandsvillie, som i Tankens Klarhed
har Inderlighedens Varme og Selvbestemmelsens Virkelighed,
er sikkert: Sandhedens Maal for Aand og Tro er
en kraftig Villie.

% --- p. 65 ---
\setcounter{footnote}{0}%
\lecture{Sjette Forelæsning.}\opage{65}

En kraftig Villie vil virke: vælge, beslutte, handle;
det vil Driftvillien, det vil Fornuftvillien, vil i høieste Forstand
Aandsvillien. Hvad vælger da Aandsvillien? Den
er i Væsen guddommelig, og vælger det Menneskelige; den
er i virksom Aabenbarelse menneskelig, og vælger det Guddommelige:
den er i sit Valg gudmenneskelig. Det er i
Ordet som Gud, at den Evige selv udvælger hvert Menneske,
og det er igjen formedelst Ordet, at Mennesket af Aandens
Kraft kan gjøre Udvælgelsen virkelig. Først da, naar den
menneskelige Villie paa een Gang vælges og vælger, vælges
af Gud til at vælge Gud, vælge sit Udspring: da brister
Egenvilliens Selvbedrag. Selvet vælger sig selv efter sin
evige Bestemmelse, vælger Frihed i ubetinget Afhængighed,
vælger at ville Eet: være Aandens Organ, vælger at ville
Guds Villie. Valget er ikke mellem Dette eller Hiint
af mange Goder; Valget er indenfor Villien selv mellem
lavere og høiere, forkrænkelig og uforkrænkelig Villie. Det
er Villielivets dybeste Krise: den gamle Villie døer, for at
fødes paa ny; Driftvillien begraves, for at opstaae forklaret
som Aandsvillie. Et saadant Valg mellem Liv og Død
er Valget mellem Tro og Vantro.

Troen, det kan ikke noksom gjentages, er mere end
en Rørelse, mere end en troskyldig Antagelse af et foreskrevet
Noget, hvorpaa der skal troes; Tro er Fordybelse
% --- p. 66 ---
\setcounter{footnote}{0}%
\opage{66}i Aand, er Hjertets fulde Hengivelse til Livets Ord; denne
Hengivelse er mere end en Indseen, mere end en Tænken,
den er en Villiesakt, en indre Gjerning, Aandsvilliens første,
oprindelige Handling. »Giv mig dit Hjerte!« det vil sige:
giv mig din Villie.

Den kraftige Villie har altid et Maal; det har Driftvillien,
det har Fornuftvillien, det har i allerhøieste Forstand
Aandsvillien. Hvilket Maal har da Aandsvillien sat?
Den vil det Timelige, men for det Eviges Skyld, det Evige
er Maal, det Timelige er Middel. Maalet er: Personlighedens
timelig betingede, men kraftig dybe, alsidig virkelige
Udfoldning af et evigt Livs uudtømmelige Indhold.
Derfor er dette Maal baade timeligt og evigt.

Driftvillien har sit Maal i det Timelige, den vil i al
Naturlighed et Nærværende; Fornuftvillien har sit Maal i
et nærværende Evigt, hvis Love blodsuge det Naturlige og
fortære det Timelige. Driftvillien vil umiddelbart det
Menneskelige, men gjør det Personlige timeligt, og bliver
umenneskelig. Fornuftvillien vil i den almene Nødvendigheds
Navn det Menneskelige, men opløser Selvet og bliver
derved umenneskelig. Det Umenneskelige, der ligger skjult
paa Bunden, saa vel under Fornuftvillien som under Driftvillien,
hidrører fra, at begge saa udelukkende ville det
Menneskelige, kun det Menneskelige. Aandsvillien derimod
har i Maalets Eenhed et dobbelt Maal; den vil i personlig
Forstand det Menneskelige, fordi den personlig vil det
Guddommelige. Den vil ikke et Guddommeligt for sig og
et Menneskeligt for sig; den vil i den sande Personligheds
hellige Navn det Gudmenneskelige. Saa omfatter dens
Maal baade Jord og Himmel, baade et Nærværende og et
Tilkommende, baade et Her og et Hisset: det dobbelsidige
Maals gudmenneskelige Eenhed er Livet selv, det
evige Liv.

At det evige Liv er Livet selv, kan sees paa Lovene.
Det forholder sig nemlig ikke saaledes, som Mange synes
at antage, at der efter Oprindelsen skulde være to Slags
Liv, et der udspringer af Timelighed, ved Siden af et
% --- p. 67 ---
\setcounter{footnote}{0}%
\opage{67}andet, som er fra det Evige. De, der holde fast paa
Forestillingen om et Liv, hvis Oprindelse, Udvikling og
Afslutning udelukkende skulde tilhøre Timeligheden, have
ikke gjort sig tilbørligt Rede for, hvad Timelighed er, og
hvad der forstaaes ved det, vi kalde Liv.

Timeligt er ingenlunde alt det, som foregaaer i Tiden;
thi enhver Udvikling, ogsaa af et evigt Indhold, maa foregaae
i Tid: Tid hører med til Evighed. Timeligt derimod
er det, som i Tiden foregaaer saaledes, at dets
Virken og Væren forudsætter Udstykning og betinges af
Sammensætning. Foranderlighed, Omskiftelighed, Forkrænkelighed,
det Enes Fremgang af det Andet og Opløsning
i det Andet, altsaa udstykket Andethed, er Kjendemærket
paa Timelighed. I denne Forstand ere alle materielle
Dannelser timelige.

Anderledes med det Evige, hvis Væsen i Alt er sig
selv ligt. Allivet er evigt, et Selvvæsen, som igjennem
Stoffet, det timelige Andet, udvikler sig selv, kun sig selv
og sit Indhold. Livet i alt Levende begynder indenfra og
holder sig under Stoffets Bearbeidelse, Omformning og
Tilegnelse i sit indre Kredsløb; i dette Kredsløb viser sig
det Evige. Der declameres om Materiens Evighed, om de
evige Stjerner, de evige Bjerge o. dsl.; det er kun Declamation.
Materiens, Stjernernes, Bjergenes Tilværelse er
en selvløs Tilværelse, en Tilværelse i Andethed, altsaa i
Timelighed. De Billioner Gange Trillioner af Aar, hvormed
man saa let kan rutte, formaae ikke at gjøre al denne
Timelighed til en virkelig Evighed. Men Allivet, som
midt i Andetheden beseirer Udstykningen, fortsætter og
bevarer sig selv, er evigt. Der er i Plantens Celler og i
Markens Blomster mere Evighed end i alle »de evige
Stjerner« tilhobe. Livsphænomenerne ere forkrænkelige, men
Livets Væsen i sit stille Kredsløb er uforkrænkeligt; det
sammensættes ikke af et Andet og løses ikke op i et
Andet: det bliver i sig selv, uforanderligt i sine Forandringer,
uforgængeligt i sit Indre, skjøndt det Ydre
forgaaer.

% --- p. 68 ---
\setcounter{footnote}{0}%
\opage{68}Man besinde sig ret paa Forskjellen mellem det Levende
og det Livløse, og man vil overbevise sig om, at
ethvert Væsen, vi ansee for et Levende, om det saa kun
lever i en enkelt Celle, bliver at ansee for et Selvvæsen,
medens Stoffet, det Materielle, med dets Kræfter og Love
kun er et Andetvæsen, er livløst, fordi det er selvløst.
Heri sees en gjennemgaaende Forskjel mellem physiske
Love og Livslove. De physiske Love ere Timelighedslove;
kun Selvudviklingslovene ere sande Evighedslove. Af blotte
Stofdele og physiske Love kan Livets Oprindelse ikke afledes:
at aflede Livet af det Livløse, Selvvæsenet af det
Selvløse er ligesaa umuligt, som at aflede Noget af Intet.
Først, idet det Timelige gaaer ind under det Evige, og de
physiske Love optages i Love for Selvudvikling, erkjende
vi Livets Love.

Hvorpaa beroer saa Adskillelsen mellem timeligt og
evigt Liv? Ikke paa Allivets Væsen, thi det er fra Evighed
af et Selvvæsen, men paa Livets særlige Fremtræden i de
levende Individer. I denne Henseende er Forskjellen mellem
Planteliv, Dyreliv og Menneskeliv af gjennemgaaende Betydning.
I Plantelivet er Selvvæsenet saa flydende, at et
bestaaende Selv ikke danner sig; i Dyrelivet faaer Selvvæsenet
et naturligt Selv, thi Dyret har en Sjæl; men
først paa Menneskelivets Trin er Spiren lagt til Udvikling
af et Selv, der er sig bevidst, et aandigt, personligt Selv.

Det Spørgsmaal, hvorledes de levende Individer kunne
fortsætte deres Selvudvikling ud over Timelighedens Grændser,
kan ingen Videnskab besvare. En saadan i det Uendelige
fortsat Udvikling af det individuelle Liv er fra de
naturlige Erfaringers Synspunkt usandsynlig, men, ifølge
Livslovene, dog en Mulighed. Det er denne Mulighed,
Troen griber, idet den opdager Selvets Evighed, opløser
den i Timelighed hildede Egenvillie og forvandler den til
Aandsvillie. Den fulde Forvisning om Individernes evige
Liv er betinget af Aandsvillie.

Det er en letfærdig Forestilling, at der i dette Nærværende
ikke skulde gives nogen personlig Vished om et
% --- p. 69 ---
\setcounter{footnote}{0}%
\opage{69}evigt Liv, men først efter Døden. Skulde vi først efter
Døden opdage, at der er et evigt Liv, kunde Opdagelsen
blive forfærdelig.

Egenvillien, hvad enten den for Øvrigt optræder i
Form af Driftvillie eller som Fornuftvillie, fornegter det
evige Liv, og gjør Brud paa dets Love. Driftvillien gjør
Individualitetens timelige Krav gjældende, men udelukker
de evige; det er et Lovbrud, et Brud paa Livsloven. Fornuftvillien
gjør Evighedskrav gjældende, men paa en selviskselvløs
Maade til Individualitetens Undergang; det er et
Lovbrud, et Brud paa Livsloven. En saadan Egenvillie
kan gjøre Væsen af Nødvendigheden, bukke og skrabe for
Naturlovenes Uforanderlighed, saa meget den vil, den gjør
dog Brud paa en uforanderlig Lov, Frihedens, Personlighedens
hellige Lov, det evige Livs Grundlov. Men at
gaae gjennem Døden ind til Evigheden som en vitterlig
Forbryder mod det evige Livs Grundlov, er dog vel en
betænkelig Sag.

Paa Aandsvilliens Trin stiller Forholdet mellem det
Timelige og det Evige sig som et Forhold i Splid og Forsoning.
Kun forsaavidt det Timelige i enhver Henseende
gjøres til Middel for det Evige, kan Aandsvillien bruge det
Timelige og trods Modstand, Hindringer og Gjenvordigheder
forlige sig dermed. Saa kan der strides og lides i det
Timelige, men naar der strides og lides for det Eviges
Skyld, erkjendes det, at Strid og Lidelse afgive Betingelser
for at udvikle Villiens Kraft og forhøie dens Styrke; med
denne Erkjendelse voxer Forvisningen om evigt Liv: det
Evige kommer bestandig nærmere, saa er der en Forsoning.
Men endskjøndt Livet i denne Forsoning er \emph{evigt},
saa er det dog endnu ikke \emph{saligt}. Saligheden er kun i
lyse Øieblikke; Spliden er ikke hævet: Modsætningen
bliver mellem Saligt og ikke Saligt. Det er Villielivets
Lov inden Timelighedens Grændser.

Anderledes, naar Timeligheden omskiftes med Evigheden,
naar \emph{Her} ombyttes med \emph{Hisset}. I Evigheden
selv ere de to Væsener, det timelige og det evige, Fjender
% --- p. 70 ---
\setcounter{footnote}{0}%
\opage{70}paa Liv og Død, uforsonlige Fjender. I Evighedens Rige
taales ingen Timelighed, thi her taales ingen Forkrænkelighed;
det følger af Sagens Natur, det ligger i Evighedsloven.
Al udvortes Timelighed, alle Mærker af naturlig
Forkrænkelighed kan Døden borttage; men det Forkrænkelige,
der er i det Indvortes, naar Selvet med sin Egenvillie
er i Forkrænkelighed, trods al sin Arbeidsvigtighed dog i
sit Selvbedrag lutter Endelighed, lutter Timelighed, lutter
Forkrænkelighed, dette Mærke kan Døden ikke borttage;
thi Døden har ingen Magt til at udslette det Personlige i
Menneskets Væsen.

Saa er da dette det evige Livs ubrødelige Lov:
»\emph{Enten --- Eller}«. Enten har Mennesket sit inderste
personlige Liv i Aandsvillien, og da kommer Evighedens
Vaar og udfolder Livet til Blomst; det er det Salige. Eller
Mennesket har sit Liv i sin Egenvillie, har i Egenvilliens
Almene villet sig selv og sin Opløsning. Saa kommer
dette Selv med sit Selvbedrag som Evighedens Fjende ind
i Evigheden. Det vil opløses, og kan dog ikke opløses,
thi det Evige i Selvet gjør sin Ret gjældende. Det har
ikke villet et evigt Liv, saa maa det lide alle Selvmodsigelsens
Kvaler; det har ikke villet Alkjærlighedens Villie,
saa møder det Magten, den dømmende Villies Almagt; det
er det Usalige. Timelighedens Modsætning af Saligt og
Ikke-Saligt kan altsaa i et Hisset forvandle sig til en
Evighedsmodsætning mellem Salige og Usalige; det er en
Dom efter Loven.

»Men«, sige den oplyste Fornufts Talsmænd, »dette
Usalige er jo netop det Inhumane. Troer Aandsvillien at
kunne kyse os med Usalighed, reent ud sagt, med Muligheden
af en evig Fordømmelse, saa have vi da her et
haandgribeligt Beviis for, hvad det er, hvortil det Hele
tilsidst peger hen. En evig Fordømmelse! Det er et nyt
Opkog af den gamle Modsigelse, en Begrebsmodsigelse,
som Videnskaben forlængst har løst, --- det Usalige, det
Barbariske, det Umenneskelige!« Vort Svar er kort. Som
Livet er, saa ere dets Love, og Lovene ere jo --- »ufor%
% --- p. 71 ---
\setcounter{footnote}{0}%
\opage{71}anderlige«. Dersom Frihedens, Personlighedens Grundlov
er barbarisk, saa er ganske vist det evige Liv barbarisk;
men saa er ogsaa al Tale om indre Selvstændighed, personligt
Maal, Kamp for Frihed, og hvad man ellers vil,
naar der tales i Alvor, meningsløs, barbarisk Tale.

Lad os holde os til det, vi kjende, det evige Livs
Love indenfor Timeligheden, og til Slutning kaste et Blik
paa Forholdet mellem \emph{det Religiøse} og \emph{det Humane}.

Vi tør da, hvad dette Forhold angaaer, ikke blive
staaende ved det Udvortes og kalde Historien til Vidne
paa, hvor ofte Religionen har villet en Vei, og Humaniteten
en anden. Spørgsmaalet er: hvad fordrer Religionen,
og hvad fordrer Humaniteten? Lade de religiøse og humane
 Fordringer sig forene eller ikke? Intet kan være mere
aandssvækkende end den i taaget Fromhed undfangede
Forestilling, at Livet i Gud og Livet i Verden skulde være
uforenelige Ting, at et kraftigt, dygtigt, virksomt Liv i den
virkelige Verden skulde være en Synd. Naar nu de, som
kalde sig »Guds Børn«, desuagtet ere nødsagede til at
»arbeide og see paa Profiten« lige saa vel som »Verdens
Børn«, saa trøste de sig med »den haarde Nødvendighed«
og »Guds Naade«, og leve hen i et gudeligt Smaavrøvl om
»Syndernes Forladelse«, som ikke fremmer det Humane, og
ingen Hjemmel har i det Religiøse. Thi hvad er det,
Religionen oprindelig kalder Synd? Det er jo den Tilstand,
hvori Mennesket befinder sig udenfor Aanden, den
Tilstand, hvori Individet er hjemfaldet til Driftvillien og
den af Gud forladte Egenvillie. Mennesket, siger Religionen,
fødes i Synd; thi det fødes under Driftvillien. Af
Tilstanden Synd udspringe de saakaldte actuelle Synder;
de gaae gjennem Følelse, Phantasie og Forestilling og modtage
den individuelle Tilregneligheds Præg, idet de optages
og fuldbyrdes af Egenvillien.

 Betragte vi disse virkelige Synder nøiere, da viser det
sig, at de lige saa vel forstyrre Livet i Verden som Livet
i Gud, ere i lige Grad ødelæggende for det Humane og
for det Religiøse. Det er en egen Art gudelig Naivetet,
% --- p. 72 ---
\setcounter{footnote}{0}%
\opage{72}naar en Sjæleryster, som rigtig vil gaae løs paa Verdens
Ugudelighed, giver sig til at dundre mod Mord, Tyveri,
Overdaadighed, Gjerrighed, Ondskab, Svig, Løsagtighed
osv., som klare Beviser paa Oprør mod Gud, Overtrædelser
af hans strenge Bud. De beængstede Samvittigheder faae
et Indtryk af, at dersom den Almægtige ikke saa bogstaveligt
havde forbudt alt Sligt i sin hellige Lov og truet
Overtræderne med Helvede, saa havde slige grove Udbrud
af menneskelig Egenvillie ikke i egenlig Forstand været
Synder. Det glemmes, at Almagtens Bud ere Kjærlighedsbud,
at det Onde forbydes, fordi det i sig selv er ondt,
og er ondt fordi det undergraver det Menneskelige. En
Gud, hvis Villie fordrer noget Andet af Mennesket end
det, som hvert Menneske, der i Tro erkjender sin evige
Bestemmelse, maa fordre af sig selv, er en Afgud; med
andre Ord: en Religion, der staaer i Strid med sand
Humanitet, er lige saa falsk som en Humanitet, der staaer
i Strid med sand Religion.

Denne Erkjendelse hænger paa det Nøieste sammen
med den rette Forstaaelse af, hvad Syndsforladelse er.
Det er en barnlig Forestilling, at Faderen i Himlen, der
er blevet vred, fordi hans strenge Bud er overtraadt, kan,
idet Forbederen træder imellem, formildes og blive god
igjen, naar Synderen kryber til Korset og beder ydmygt
om Forladelse. Det Barnlige i en saadan Forestilling bør
vi ære og agte, thi det er evig vist, at Gudsforholdet fra
Menneskets Side vedblivende maa være et barnligt Forhold.
Men det Barnlige er væsenlig forskjelligt fra det Barnagtige.
Og barnagtigt bliver Forholdet saa vel hos de
»søde« Christne som hos de »sure«, naar Syndsforladelsens
evige Grundlov miskjendes.

De Glade, der ere saa rødmussede, fordi de midt i
»Verdensvrøvlet« kunne gjøre saa omtrent, hvad de ville,
og dog altid være forvissede om »Syndernes Forladelse«,
fordi den staaer i den tredie Troens Artikel, ere ikke
barnlige, men barnagtige. Istedenfor at maale Synden
med Selvmodsigelsens lige saa humane som religiøse
% --- p. 73 ---
\setcounter{footnote}{0}%
\opage{73}Maalestok, og kjende Angerens Alvor og Syndsforladelsens
Virkelighed paa den gjenopvakte Aandsvillies fornyede
Spændkraft, skyde de Loven tilside, og voxe og groe i den
Indbildning, at en gammel Psalme er vel saa god som
en ny Villie.

De sure Christne, der holde paa Tvangslovene, have
Besværlighed nok, ikke blot med deres egne, men ogsaa,
og det fornemmelig, med Andres Synder: intet Under,
at de ere saa blege. Naar det store Synderegister med
Hor og Mord og Luxus og Overdaadighed og Svig og
Tyveri osv. oprulles, undgaaer det deres Opmærksomhed,
at Sjæledommeren staaer med den ene Fod i Himlen,
men med den anden paa Raadstuen; det undgaaer dem,
at de store Synder paa Regnebrættet i Grunden ere lige
saa inhumane som irreligiøse, at Politiet stadig er paa
Spor efter de groveste, uagtet Politiet jo dog ikke er en
religiøs, men en slet og ret human Institution: denne
storladne Ivren mod »Synderne« og »Synderne« og »Synderne«
er, naar den sees i sit Forhold til Aand, mindre
alvorlig end barnagtig.

Jo mere den religiøse Erkjendelse klarer sig, desto
tydeligere vil det indsees, at Livet i Verden maa være en
væsenlig Betingelse for Livet i Gud, saa sandt det religiøse
Liv i høieste, fyldigste Maal er et Villieliv. Eet er \emph{fornødent};
men for dette Enes Skyld er et Verdensindhold
af Følelser, Forestillinger, Interesser og Arbeider fornødent.
Synden kan nu antage to Former. Enten mangler Selvet
det ene Fornødne, Aandskjærligheden med Tro og Haab
er udeladt, og Egenvillien kaster sig graadig over det
Mangfoldige: det er Synden i Form af Verdslighed;
eller Selvet forbruger al sin Villiekraft paa at skille sig
af med det virkelige Verdensindhold og forbrænde i Gud,
det være nu ad Endelighedens Vei ved vilkaarlige Forsagelsers
gode Gjerninger, eller ad Uendelighedens Vei ved
mystisk at styrte sig i den unævnelige Alkjærligheds
Lyshav: det er Synden i Form af misforstaaet Gudelighed.
Livet i Gud betyder, at Mennesket ubetinget
% --- p. 74 ---
\setcounter{footnote}{0}%
\opage{74}har opgivet sin Egenvillie og erkjendt sin absolute Afhængighed
af Guds Villie. Livet i Verden betyder, at
den absolute Afhængighed begrunder en virkelig Frihed,
at Aandsvillien gjør Mennesket selvstændigt. Aanden
beviser just sin Personligheds Kraft deri, at Mennesket
med kraftig Villie og sikker Haand tager fat paa det
virkelige Livs Arbeider.

Middelalderens naive Adskillelse mellem \emph{gudelige}
og \emph{verdslige} Gjerninger kan hverken bestaae for Religionens
eller for Humanitetens Dom. Alle sædelige Handlinger
ere i lige Grad religiøse og humane, saa sandt de
udgaae fra Aandsvillien.

Ingen Gjerning i Verden, der staaer i noget Forhold
til det personlige Liv, er Aanden for ringe; og ingen
udvortes endelig Gjerning, selv om den overbød den største
i en Loyola's himmelstormende Askese, er stor nok til
at udtrykke Aanden som Aand.

Der er i alle Gjerninger, de være »gudelige« eller
»verdslige«, Synd, naar Aandsvillien mangler; men i alle
Gjerninger, de kaldes nu verdslige eller gudelige, ophæves
det Syndige, forsaavidt Aandsvillien træder i Kraft.

Der er ikke tre Slags Gjerninger: nogle, vi skulle
gjøre for Guds Skyld, andre vi skulle gjøre for Næstens
Skyld foruden dem, vi skulle gjøre for vor egen Skyld.
Elsk Gud over alle Ting og din Næste som dig
selv! Disse Lovens store Bud ere eet Bud. Det er
umuligt at elske sig selv i Sandhed uden at elske Gud,
og det er umuligt at elske Gud uden at elske sin Næste.
Det Religiøse er, at Kjærlighed til Gud staaer øverst; det
Humane er, at Kjærlighed til os selv staaer uadskillelig
fra Kjærlighed til Næsten.

Det gamle ufrugtbare Skolespørgsmaal, hvorledes vi
med Sikkerhed skulle udfinde, hvad Guds Villie er med
Hensyn paa Livets enkelte Gjerninger i de forskjellige
mødende Tilfælde, falder fra dette Synspunkt bort af sig
selv. Tilværelsen stiller Mennesket en uoverskuelig Mængde
af høist forskjellige Opgaver. Ingen af disse Opgavers
% --- p. 75 ---
\setcounter{footnote}{0}%
\opage{75}Løsning er overladt til Menneskets selviske Forgodtbefindende.
Det Religiøse fordrer, at de alle, umiddelbart eller
middelbart henføres --- i Troen --- til Guds Villie. Men
ingen af dem alle bliver i sin umiddelbare Enkelthed bogstavelig
bestemt ved en af Gud dicteret Befaling; de skulle
alle opfattes med Aandens Frihed. Det Humane skeer sin
Ret, idet den menneskelige Villie som Aandsvillie er i
hvert mødende Tilfælde opfordret til at handle ud af
Personlighedens Indre.

Grundforholdet mellem det Religiøse og det Humane
lader sig, ifølge vort Foregaaende, nøiagtig bestemme ved
Forholdet mellem Driftvillie, Fornuftvillie og Aandsvillie.
Fornuftvillien betegner, som vi have seet, et Overgangstrin,
og lader sig ikke gjennemføre; den svæver bestandig
mellem Religiøst og Irreligiøst, mellem Humant og Inhumant.
Driftvillien vil, idet den begynder at blive fornuftig, uden
videre ind paa det Humane. Om den end bliver paavirket
af Aand, opgiver den derfor ikke sig selv, men gaaer paa
Accord med det Religiøse, og søger paa bedste Maade at
drage Fordeel af Religionen for at bevare og fremme Humaniteten.
Hertil svarer den ofte noget krinkelkrogsagtige
Brug, som politiserende Samfund, af Klogskabshensyn, have
vænnet sig til at gjøre af de religiøse Institutioner.
Aandsvillien derimod staaer i et dobbelt Forhold til det
Humane, et fjendtligt og et forsonligt. Den staaer i Strid
med Drifthumaniteten, thi den staaer i Strid med den
naturlige Driftvillie. Den stiller sig afgjørende paa Religionens
Side, retter sine Angreb mod den aandsfornegtende
Humanitet, nedbryder dens Altere og brænder dens Idoler.
Saa seer det ud, som om den var barbarisk. Men disse
Angreb paa falsk Humanitet ere i Principet til Bedste for
det i Sandhed Humane; just derved, at det Humane
ubetinget gaaer ind under det Religiøse, bliver det gjenfødt
til et høiere Liv: Aandsvillien griber ind i det Verdslige,
planter sit Frihedsbanner og gaaer i Kamp for Humanitetens
Ret.

% --- p. 76 ---
\setcounter{footnote}{0}%
\opage{76}Det tvesidige Forhold mellem Aandsvillien og Humaniteten
lader sig eftervise i alle Samfundets Kredse, i
Familien, i det Borgerlige, i de Almeenhedens Formaal,
for hvilke der arbeides med Ideer.

Ægteskabet er Familielivets sædelige Grundlag.
Hvorfra stammer det? hvilken Gyldighed har det? er det
opløseligt eller uopløseligt? De, som leve med den usminkede
Drifthumanitet, maae naturligviis erklære Ægteskabet
for et naturligt Produkt af Driftvillien. Det Humane er
da, at Ægteskabet betinges af gjensidig Kjærlighed; de
umiddelbare Hjerters umiddelbare Hengivelse med Offer
og Accidenser er Ægteskabets Lykke. Hændes det nu,
at Kjærlighedsilden gaaer ud, og det kan saa let hændes,
da det jo er Driftkjærlighed, saa maa Ægteskabet hæves,
for ikke at blive ulykkeligt. Have Individerne ikke uden
videre Ret til at hæve det selv, da er det naturligviis en brutal
Inhumanitet af et slavisk Bestaaende, som tør understaae
sig at sætte Skranker for den frie Kjærlighed i »det frie
Kjød«. At Aandsvillien ikke kan forlige sig med en saadan
Humanitet, er begribeligt. Men det Samfund, der er
blevet staaende ved at faae Drifthumaniteten sminket op med
nogle sædelige Skjønhedspletter, uden at ville gaae ind
paa Religionens ubetingede Fordring, tilfredsstiller heller
ikke Aandsvillien. Fra dette Standpunkt faaer man ud,
at Ægteskabet ganske vist, efter den guddommelige Lov,
er uopløseligt, men dog, naar der indhentes speciel Bevilling,
være sig fra Paven eller Cancelliet, alligevel opløseligt
af Hensyn til det Humane. Imod en saadan halvreligiøs,
forhutlet Humanitet have alvorlige Bodsprædikanter
til forskjellige Tider hævet deres myndige Røst: »det
guddommelige Bud maa høres, erkjendes og lydes, om saa
Alverdens Humanitet skal gaae til Grunde derved!« Saa
seer det ud, som om Religionen var grusom og foragtede
det Humane. Og dog er det en uhyre Misforstaaelse.

Det Religiøse vil herske absolut, ubetinget; men ikke
ved udvortes Lovbud. Det udvortes Lovbud, hvad enten
det bygges paa Skriftsteder eller paa Cancellirescripter eller
% --- p. 77 ---
\setcounter{footnote}{0}%
\opage{77}paa begge Dele, maa ifølge sin Natur altid holde sig
indenfor endelige Betingelser, og kan i sin Udvorteshed
aldrig gjøre Fordring paa at gjælde for absolut. Det
Absolute ligger i, at det guddommelige Bud i Aandens
Navn er et Kjærlighedsbud, et Personlighedens, Selvstændighedens
og Frihedens Bud. Den religiøse Forstaaelse
er, at Budet i al sin Strenghed ikke vil paalægge Individerne
udvortes Tvang, men overbevise dem om Aandsvilliens
Krav og fange dem i deres egen Samvittighed. Kommer
en religiøs Individualitet efter alvorlig Selvprøvelse for Gud
til den ydmygende Vished, at det indgaaede Ægteskab er,
af hvilkesomhelst Grunde, blevet et tomt Skin, et Trældomsaag,
en Snare, da tør en saadan Personlighed, være
sig Kvinde eller Mand, træde selvstændig op, bryde sit
Aag og gaae sin egen Vei trods Tungernes Travlhed og
et retfærdigt Nabolavs velbekjendte Pharisæisme: det Mod
skal Troen give. Det er Religionens Humanitet, og
sandere Humanitet kan ikke gives.

Det borgerlige Samfunds sædelige Grundlag er Ret
til Eiendom; hvor Eiendomsretten ikke agtes, kan ingen
Samfundsorden bestaae. Den driftige Borger vinder Indflydelse
i Forhold til som han udvider sine Besiddelser.
Den borgerlig juridiske Personlighed fremtræder i Form
af Eiendom: jo større Eiendom, desto mere Personlighed.
Det er naturlig Humanitet. Driftvillien har sit vigtigste
Mellemværende i Maalforhold af Mit og Dit. Saa danner
der sig et humant Samfund med humane Formaal. Kaldes
nu det Religiøse til Hjælp for at indvie og velsigne al
denne Humanitet, faaer den personlige Eiendom naturligviis
høie Tanker om sig selv. Eiendommen bliver Capital,
og Capitalens Magt er uimodstaaelig. Capitalen bugner
under Velsignelsen; den voxer og formerer sig, gjør sig
Jorden underdanig og regjerer over Fisken i Havet og
Majestæten paa Thronen, over Arbeideren i Gruben
og over hvert ubemidlet Dyr, som kryber paa Jorden.
For at sikkre sig alle Velsignelsens Frugter, bliver Drifthumaniteten
gudelig: den rige Mand gaaer i Kirke om
% --- p. 78 ---
\setcounter{footnote}{0}%
\opage{78}Søndagen, giver Summer til det Offenlige osv.; paa
Driftvilliens Standpunkt altsammen hæderligt, smukt og
humant. Men hvad siger Religionen? Bort med det altsammen!
»Sælg alt dit Gods og giv det til de Fattige ;«
»det er lettere, at en Kameel gaaer igjennem et Naaleøie
end at en Rig kommer ind i Guds Rige!« --- Hvor forargeligt.
Hører I, sige de Fornuftige, hvor fanatisk
eensidig Religionen er, hvor lidt den forstaaer sig paa det
i Virkeligheden Humane!

For at tage det Forargelige bort og berolige Sindene
kommer en blødsøden Sjæletrøster med den ligesaa videnskabelige
som humane Fortolkning, at det ingenlunde er
Religionens Mening, at vi bogstavelig skulle bortgive alt
vort Gods: den vil kun vække vor Samvittighed og
sætte de Gjerrige Skræk i Blodet. Den rette Mening er,
at vi skulle betænke de Fattige, give Noget bort af Godset,
en god Deel, især til gudeligt Brug, men ikke Alt; det
Meste skulle vi selv beholde, at vi ret maae lære at bestyre
det til Gavn og Glæde for os selv og vore Arvinger.
Men hvad siger Religionen? Bort med det Hele! »Sælg,
hvad Du haver, og giv de Fattige det.« Det er Veien
til Fuldkommenhed, Forsagelsens Vei til »Himmeriges
Rige«. Saaledes blev Budet forstaaet i hine Dage, da
Ørkenens Eremiter og Middelalderens strenge Munke endnu
ikke havde lært at fuske i det religiøst Humane.

Dog, indvender man, »det store Fremskridt, Nutiden
har gjort, bestaaer jo netop i, at den vil forlige Religion
med Humanitet.« Tilvisse! Men Forliget skal dog vel
ikke grunde sig paa en underfundig Svig. Misforstaaelsen
hos hine Gamle var ingenlunde den, at de toge den religiøse
Fordring for ubetinget. Paa den bogstavelige Forstaaelse
staaer en hellig Antonius, en Benedict fra Nursia,
en Francesco af Assisi uendelig høiere end alle den
moderniserede Christendoms Mæglingsethikere tilhobe.

Asketernes Misforstaaelse var ingenlunde den, at
de toge Forsagelsens Bud for strengt, men tvertimod den,
at de ikke toge det strengt nok. De opfattede Verdens%
% --- p. 79 ---
\setcounter{footnote}{0}%
\opage{79}forsagelsens Bud som et udvortes Bud, et til Fuldkommenheds
Opnaaelse vilkaarlig foreskrevet Tvangsbud. Vel
var dette Bud dem ogsaa et Hjertets Bud, et Kjærlighedsbud
til Befrielse og Salighed; men Kjærlighedens
Bud blev ikke opfattet i Kjærlighedens Aand, thi det
blev ikke opfattet i den personlige Friheds Aand. Saa
seirede Ufriheden og der fremkom indenfor Religionen to
Slags Moral, en høiere for de Fuldkomne og en lavere
for de Ufuldkomne. Det var dette Misgreb, der gjorde
den middelalderlige Religiøsitet inhuman.

Først naar Forsagelsens Bud bliver tilegnet i Tro
som et indvortes Bud, et Frihedens, Kjærlighedens og
Personlighedens Bud i al sin Strenghed, bliver Strengheden
igjen Mildhed, og det Humane gjenfødes i det Religiøse.
Det seer saa strengt ud, at Mennesket skal i udvortes
Forstand skille sig ved alt sit Gods; Driftvillien seer ingen
Bagdør, hvorigjennem den kan slippe bort fra en saa nærgaaende
Fordring. Naar det Udvortes derimod omsættes
i det Indvortes; naar Budet vender sig og siger: behold,
hvad du har, i det Udvortes; men slip det i det Indvortes,
vær med al din Eiendom som den, der ikke eier en Hvid,
saa, mener man, er Religionen saare mild: Driftvillien øiner
Bagdøren, smiler i Skæget, og har sine egne Tanker. Og dog
er det saa, at det hellige Kjærlighedsbud, forstaaet i Aand
og Sandhed som et Frihedens Bud til personlig Selvudvikling,
er i al sin Mildhed langt strengere end et hvilketsomhelst
asketisk Tvangsbud. Det er Aandsvilliens Kraft,
der skal prøves i Humaniteten.

Dersom det religiøse Princip for sand Humanitet
ikke bliver anerkjendt, vil Striden om Mit og Dit udvikle
sig til en barbarisk Strid, og i en barbarisk Strid vil Seiren
tilsidst være paa de elementære Kræfters Side. Bliver
Humanitetsprincipet derimod tilfulde anerkjendt, ikke mindre
paa de Besiddendes end paa de Ikke-Besiddendes Side,
vil Verdenskampen gaae sin Gang, men Humaniteten,
støttet paa Aandsvilliens Kraft, vil seire. 

% --- p. 80 ---
\setcounter{footnote}{0}%
\opage{80}Kaste vi et Blik paa de forskjellige Arbeider, som
udføres i Ideernes Tjeneste og ere betingede af særlig
Begavelse, saa viser sig Modsætning mellem Idee og
Aand. Hvor stor denne Modsætning er, kan sees paa
Grækere og Jøder. Grækerne have en overstrømmende
Rigdom af frugtbare Ideer, men, hvad vi i egenlig Forstand
maae kalde Aand, have de ikke. Jøderne staae
paa Aandens Trin under den lovgivende Jehovas Villie,
men have en paafaldende Mangel paa productive Ideer.
Jøderne repræsentere det Religiøse, Grækerne det Humane.
Man vil desaarsag finde det naturligt, at Christendommen
under sin store Indtrædelseskamp i Verden begyndte med
at bryde Staven over alle Hedenskabets Genieværker; det
var i de Dannedes Øine et talende Beviis paa fanatisk
Inhumanitet. Man vil endvidere finde det naturligt, at
den Humanitet, som en religionsfornegtende Nutid vil
grundlægge paa det moderne Dogme om Drift- og Fornuftlivets
absolute Eenhed, vender sig med Ringeagt bort fra
»den semitiske Races« stivsindede Borneerthed og søger
sine Forbilleder under den græske Himmel. Man vil
endelig finde det i sin Orden, at de, som ikke her tør
see det Fjendtlige i Forholdet mellem en aabenbaret Religion
og en hedensk Humanitet, fordi de ikke tør vedgaae
for sig selv, at to absolut ueensartede Principer maae kæmpe
paa Liv og Død, indtil enten det ene eller det andet faaer
Magten, hutle sig igjennem ved en Mægling efter deres
eget Hoved.

Paa Halvreligiøsitetens lige saa halvhumane som principløse
Mæglingstrin møde vi da igjen den middelalderlige
Adskillelse mellem Gudeligt og Verdsligt, men med et
moderniseret Tilsnit. »Alle Humanitetens Ideer ere berettigede,
forstaaer sig; Talent og Genie ere gode Gaver,
og alle gode Gaver skulle bruges.« Men nu indtræder den
Vanskelighed, at de allerfleste humane Ideer ere af Naturen
rene Hedninger, og det gaaer dog ikke vel an i vor Tid
at døbe Ideerne, som man i Middelalderen døbte Klokkerne.
Man maa da gjøre behørig Forskjel mellem dem paa den
% --- p. 81 ---
\setcounter{footnote}{0}%
\opage{81}ene Side, der ikke lade sig henføre under det Religiøse ---
Physikens f. Ex. og Komikens --- og dem paa den anden
Side, som deels umiddelbart deels middelbart kunne bruges
i Religionens Tjeneste. Det første Slags maa man lade
gaae paa deres egen Haand i al Verdslighed og kun sørge
for, at de holde sig inden visse Grændser. Det sidste
Slags maa man sætte Priis paa, thi de ere Mæglere mellem
det Religiøse og det Humane, efter som de bruges
mere eller mindre gudeligt. Det skulde da, hvad Anvendelse
af Talent og Genie angaaer, ifølge denne Synsmaade,
være mere gudeligt at bygge Kirker end at bygge Raadstuer,
mere gudeligt at spille Orgel end at spille Violin,
mere gudeligt at male Helgene end at male Køer. Naar
man nu desuagtet ikke aabenlyst vil vedkjende sig en
saadan Synsmaade, da er det nok fordi man nødig vil vedkjende
sig den Forblanding, som ligger til Grund for
Synsmaaden; thi det er netop Forblandingen af, hvad
der vedkommer Ideen, og hvad der beroer paa Villien,
som forvirrer Opfattelsen saavel af det Religiøse som af
det Humane.

Ideeformens Indhold er altid fra en eller anden Side
et Verdensindhold; Gudsideen selv gjør herfra ingen
Undtagelse. --- Ingen Idee er i og for sig enten gudelig
eller ugudelig; heri ligger, at alle Ideer uden Undtagelse
kunne henføres under Aandsvillien, forsaavidt Aandsvillien
ingenlunde vil kue, men tvertimod nære, rense og beaande
den Driftvillie, der føder dem til Verden. Enhver Idee
har et Maal i sig selv, og enhver Idee er et Middel for
Aanden. At negte Idee-Udviklingen sin naturlige Ret,
er inhumant; men at sætte det absolute Maal i nogensomhelst
Idee og forgude Ideen, hvad enten det nu er en
Videnskabens eller en Konstens Idee, Ideen af et Statssystem
eller Ideen af et Kirkesystem, er irreligiøst.

Hvor Aandsvillien mangler, blive Ideerne selviske;
den Spænding og Splid, som opstaaer imellem dem, sætter
Lidenskaberne i Bevægelse og giver dem Tøilen. Et bølgende
Hav af selvisk-ideelle Bestræbelser kan sættes i Oprør;
% --- p. 82 ---
\setcounter{footnote}{0}%
\opage{82}Stormen kan stige til en saadan Grad, at et rigt begavet
Folk gaaer under i Striden. Vi have i Grækerne et
tragisk Exempel.

Den naturlige Humanitet kan ikke bære sig selv,
den maa bæres af Aanden. De mange Ideeconflicter ere
Opgaver for Villien; de danne sig for Villiens Skyld, de
opfordre den til Anstrengelse, de kræve dens Afgjørelse.
I en stadig Kamp mod selviske Hindringer have vi
Betingelserne for den personlige Udvikling af Aandsvilliens
Kraft.

Paa disse Tilværelsens Grundlove vil det i Troen
kunne sees, at der maa være en inderlig Forbindelse
mellem det Nærværende og det Tilkommende, mellem
Her og Hisset. Denne Forbindelse kan ikke begrundes
ved udvortes Betragtning af »de Dødes Opstandelse«.
Den store Umage, som christelige Dogmatikere, især de
ældre, have gjort sig med en naturligviis »videnskabelig«
Udmaling af den tilkommende Herlighed, hvorledes de
Døde nu egenlig skulle opstaae, i hvilken Alder, i hvilken
Tilstand, eenbenede eller tobenede, med Paryk eller uden
Paryk, er spildt Uleilighed % PRINTED AS IS: no stop
Det hele Spørgsmaal viser
tilbage til Inderligheden, Personligheden, Aandsvillien,
Friheden.

Det Mangfoldige er til for det Enes Skyld, det
Timelige er til for det Eviges Skyld: det evige Liv er et
indre Liv, der nærer sig af Erindring. Alt, hvad der
opleves i det Udvortes, efterlader et Spor i det Indvortes,
det vil sige: det \emph{inderliggjøres}, det erindres. Jo
kraftigere Villien er under Oplevelsen, desto dybere graves
Erindringens Spor. At disse Spor nogensinde kunne udslettes,
er en Indbildning, et sandseligt Selvbedrag, der
hidrører fra, at Organerne sløves. I Sjælens Liv, i Personlighedens
Grund udslettes de aldrig: Livets Love vise
alle hen til een og samme Grundlov, Erindringens Lov.
Meget siges og Meget høres, Meget huskes og Meget
glemmes i dette Timelige, men Aanden lytter og gjemmer;
den glemmer Intet: i Evigheden glemmes Intet. Hu%
% --- p. 83 ---
\setcounter{footnote}{0}%
\opage{83}kommelsen kan svækkes, Forstanden omtaages, og Bevidsthedens
Lys udslukkes; men dybere end alt dette og
til Grund for alt dette ligger Selvets Natur forvaret og
skjult i Aandens Erindring. Af Erindringens Nat gaaer
Alt frem igjen, og det dages paany: Selvbevidsthed,
Forstand og Hukommelse vaagner op med alt det Oplevede
og først og sidst med alt det, som i Alvor er villet.
Hvad der i dette Timelige er levet og nydt, hvorfor her
er arbeidet, lidt og stridt, hvad Villien har naaet, og
hvad den inderlig har tilstræbt, men ikke opnaaet, af alt
dette gaaer Intet tabt; trods sandselig Glemsomhed bliver
Intet glemt, det bliver altsammen gjemt, begravet og gjemt
i Erindringens Skjød til fornyet Gjennemsyn.

Vort Æmne er ikke udtømt; det er i sig selv
uudtømmeligt. Langt fra den Indbildning at kunne udtømme
det Uudtømmelige, har jeg i disse Foredrag indskrænket
mig til at belyse nogle Træk af Villielivets store
Billede for derved at yde et, om end ufuldkomment, Bidrag,
med det Ønske, at Enhver, som vil tage sine Livserfaringer
til Hjælp, maa i Alvor selv blive opmærksom paa, hvad
Aandsvillie er, og klare sig Betingelserne for en kraftig
Villie.

\end{document}
